% !TeX root = ../../Book3.tex
% 纲要标题：## 第22章　流、切片与 Federer–Fleming 紧性
% 纲要位置：几何学计划纲要.md，第 3007 行。
% 正文状态：初稿；尚未正式审核。

% 纲要内容开始
% 
% > 本章目标与证明范围：从测试形式的对偶构造流，区分一般、正规、可整流和整数流。变形定理、整数性闭包和紧性均为正文证明任务；仅用弱紧性不能完成整数 Plateau。
% 
% 纲要内容结束

\chapter{流、切片与 Federer–Fleming 紧性}
\label{g3:ch:22}
\BookChapterNumber{g3:ch:22}

\begin{chapterabstract}
	流把定向几何对象对微分形式的积分推广为连续线性泛函。它能够记录重数、边界以及取向抵消，也允许讨论不光滑对象的极限。本章从质量测度与基本运算出发，区分一般流、正规流、可整流流和整数流；建立 Lipschitz 切片、锥构造及平坦范数的对偶描述。随后通过立方网格上的径向投影和面内 BV 密度修正证明变形定理，再以零维切片的整数信息、统一 BV 控制和 Lipschitz 图覆盖证明整数流闭包。最终得到共同紧支与统一质量、边界质量界下的 Federer--Fleming 紧性定理，并说明局部版本如何容许质量逃逸。
\end{chapterabstract}

\begin{learninggoals}
	\begin{itemize}
		\item 从测试形式与余质量范数定义流及其质量，核对边界、推前、乘积和同伦公式的符号。
		\item 根据切向结构、重数与边界质量区分四类流，并把有限周长集合写成顶维整数流。
		\item 用截断后的新增边界定义切片，计算其取向、质量积分和边界相容关系。
		\item 构造锥填充，使用平坦范数比较低质量误差与低质量填充，辨别三种收敛。
		\item 完成径向选点、奇性删管和 BV 中位数修正，证明统一网格变形及整数闭流填充估计。
		\item 由原子切片的 BV 控制构造 Lipschitz 图，恢复极限的切向结构与整数重数。
		\item 运用闭包和紧性定理保留指定边界，说明局部紧性与全空间平坦紧性的条件差别。
	\end{itemize}
\end{learninggoals}

把一张很薄的矩形的边界按逆时针方向走一周，上下两条长边的方向相反。矩形高度趋于零时，两条长边都靠近同一个线段，周长却不趋于零。若拿一形式沿边界积分，上下两部分会逐渐抵消，积分反而趋于零。这个例子说明：一个有向边界的极限，既要容许位置变化，也要保留相消的代数规则。

\ProseParagraphBreak
第21章用示性函数的分布导数处理了有限周长集合。那是一种余维一的边界模型；对于空间中的曲线、带重数的曲面或更高余维对象，微分形式提供了统一的测试工具。边界可以由 \(\partial T(\eta)=T(\mathrm d\eta)\) 定义，因而连续地传到弱极限。真正的困难是极限的形状：质量与边界质量有界，为什么能保证整数重数的曲面仍然收敛到同类对象？

\ProseParagraphBreak
我们先把流变形到细立方网格。流与网格模型之差由一个小填充及一个小的边界误差控制，这给出平坦紧性。整数性还需另一条论证：沿坐标切片后得到带整数权的点；正规性使这些点的测试值满足统一 BV 估计，于是能按参数组织成 Lipschitz 图。边界质量排除非切向的极限向量，面积公式再从纤维中的整数权读出图上的整数重数。这条证明把可整流性与整数性分别落实。

\ProseParagraphBreak
本章所称 Federer--Fleming\footnote{Herbert Federer（1920—2010），奥地利出生的美国数学家，主要研究几何测度论。Wendell Helms Fleming（1928—2023），美国数学家，与 Federer 合作建立了整数流的基础理论。}紧性，是固定紧支及统一质量、边界质量界下的整数流紧性。第20章的面积、余面积与可整流密度结果，以及第21章的 BV 逼近、紧性和约化边界结构，是本章的主要几何分析基础。微分形式的外微分、拉回和光滑 Stokes 公式在使用处恢复约定；闭包论证还用到本卷已证的有界泛函延拓及测度表示。

% !TeX root = ../../Book3.tex
% 纲要标题：### 22.1 流及其几何运算
% 纲要位置：几何学计划纲要.md，第 3011 行。
% 正文状态：初稿；尚未正式审核。

% 纲要内容开始
% 
% * 测试形式、连续性和边界
% * 质量、局部质量与支集
% * 推前、乘积和同伦公式
% * 光滑定向子流形及抵消现象
% 
% 纲要内容结束

\section{流及其几何运算}
\label{g3:sec:2201}

一条定向曲线对一形式的积分，同时记录曲线的位置、方向和重数。曲线发生自交时，积分仍有意义；两段重合但方向相反的曲线则相互抵消。流把这种积分配对作为基本对象。以下先在开集 \(U\subset\mathbb R^n\) 上工作；微分形式的外微分和拉回沿用光滑几何中的约定。基本构造可对照\cite[第25--26节，第129--144页]{Simon1983}；本节质量采用余质量的对偶范数，与该书第26节采用的外代数 Euclid 范数有所区别，两者在简单向量上相同。

\subsection{测试形式、连续性和边界}
\label{g3:subsec:2201:01}
记 \(\mathcal D^k(U)=C_c^\infty(U;\Lambda^k(\mathbb R^n)^*)\)。若
\(\omega_j=\sum_I\omega_{j,I}\,\mathrm dx^I\)，则测试形式收敛 \(\omega_j\to0\) 的含义是：全部支集包含于同一个紧集 \(K\Subset U\)，并且每个系数的每阶导数都一致趋于零。这里 \(I\) 遍历严格递增的 \(k\) 重指标。

\begin{definition}\label{g3:def:current}
	设 \(U\subset\mathbb R^n\) 开，\(0\le k\le n\)。记 \(\mathcal D^k(U)=C_c^\infty(U;\Lambda^k(\mathbb R^n)^*)\)，带固定紧支上各阶系数导数一致收敛的测试形式拓扑。\term[liu]{流}[current]是测试形式空间 \(\mathcal D^k(U)\) 上的连续线性泛函，称为 \(k\) 维流，记其空间为 \(\mathcal D_k(U)\)。等价地，对每个紧集 \(K\Subset U\)，存在整数 \(r\ge0\) 与常数 \(C\)，使
	\[
		|T(\omega)|\le C\max_{I,\,|\alpha|\le r}\|\partial^\alpha\omega_I\|_\infty
		\quad(\operatorname{supp}\omega\subset K).
	\]
	若 \(k\ge1\)，定义\term[liu de bianjie]{流的边界}[boundary of a current]为
	\[
		\partial T(\eta)=T(\mathrm d\eta),\qquad \eta\in\mathcal D^{k-1}(U).
	\]
	零维流的边界规定为零；当 \(k<0\) 或 \(k>n\) 时，规定 \(\mathcal D_k(U)=\{0\}\)。流列的弱收敛 \(T_j\rightharpoonup T\) 指对每个测试形式都有 \(T_j(\omega)\to T(\omega)\)。
\end{definition}
\symidx{\(\mathcal D^k(U),\mathcal D_k(U)\)：测试 \(k\) 次形式与 \(k\) 维流空间}
\symidx{\(\partial T\)：流 \(T\) 的边界}

上述半范数判别可以直接从连续性得到。若在某个 \(K\) 上不存在这样的 \(r,C\)，逐个选取支于 \(K\) 的形式，使前 \(j\) 阶导数均不超过 \(1/j\)，而泛函值绝对值大于一，就与测试形式连续性矛盾；反向由所给估计立即得到。外微分连续，故 \(\partial T\) 仍是流。由 \(\mathrm d^2=0\)，
\[
	\partial^2T=0,\qquad
	T_j\rightharpoonup T\ \Longrightarrow\ \partial T_j\rightharpoonup\partial T.
\]
边界运算因而把几何边界与代数上的链复形联系起来。

零维流正是分布。例如点流 \(\delta_a(\varphi)=\varphi(a)\)，以及 \(\partial_{x_i}\delta_a(\varphi)=-\partial_i\varphi(a)\)，都是零维流。它们的流边界均为零；这里的分布偏导与流边界是不同运算。后者降低流的维数。

\begin{example}\label{g3:exm:oriented-segment-current}
	设 \(a,b\in\mathbb R^n\)，令 \(\gamma(t)=a+t(b-a)\)。线段流定义为
	\[
		\lBrack a,b\rBrack(\omega)=\int_0^1\langle\omega(\gamma(t)),b-a\rangle\,\mathrm dt.
	\]
	微积分基本定理给出
	\[
		\partial\lBrack a,b\rBrack=\delta_b-\delta_a.
	\]
	交换端点改变整个流的符号。相邻线段相加时，公共端点在边界中抵消；曲线是否自交并不影响这项计算。
\end{example}

\subsection{质量、局部质量与支集}
\label{g3:subsec:2201:02}
外积 \(v_1\wedge\cdots\wedge v_k\) 称为简单 \(k\) 向量，其 Euclid 长度等于这些向量张成的平行多面体体积。对一个 \(k\) 余向量 \(\omega\)，定义\term[yu zhiliang]{余质量}[comass]
\[
	\|\omega\|_*=\sup\{|\langle\omega,\xi\rangle|:\xi\text{ 简单},\ |\xi|=1\}.
\]
其对偶是 \(k\) 向量的质量范数
\[
	|v|_{\mathrm m}=\sup_{\|\omega\|_*\le1}|\langle\omega,v\rangle|.
\]
简单向量满足 \(|v|_{\mathrm m}=|v|\)：上界来自余质量定义，下界由其定向平面上的单位体积形式取得。一般向量不必简单；例如在 \(\mathbb R^4\) 中，\(|e_1\wedge e_2+e_3\wedge e_4|_{\mathrm m}=2\)，而 Euclid 范数为 \(\sqrt2\)。上界由三角不等式得到，下界用余质量为一的 \(\mathrm dx^1\wedge\mathrm dx^2+\mathrm dx^3\wedge\mathrm dx^4\) 检验。

\begin{definition}\label{g3:def:current-mass}
	设 \(U\subset\mathbb R^n\) 开，\(0\le k\le n\)，\(T\in\mathcal D_k(U)\)。下式 \(\mathcal D^k(V)\) 表示紧支光滑 \(k\) 形式，\(\|\cdot\|_*\) 表示形式在单位简单 \(k\) 向量上的上确界，即余质量范数。对开集 \(V\subset U\)，定义\term[liu de zhiliang]{流的质量}[mass of a current]
	\[
		\mathbf M_V(T)=\sup\{T(\omega):\omega\in\mathcal D^k(V),\ \sup_x\|\omega(x)\|_*\le1\}.
	\]
	记 \(\mathbf M(T)=\mathbf M_U(T)\)。若每个 \(V\Subset U\) 上此数有限，称 \(T\) 具有局部有限质量。\(T\) 的支集 \(\operatorname{spt}T\) 是所有使 \(T\) 消失的开集之并的补集，其中在开集上消失指其内每个测试形式的配对均为零。
\end{definition}
\symidx{\(\mathbf M(T),\mathbf M_V(T)\)：流的总质量与开集 \(V\) 中的质量}

\begin{proposition}\label{g3:prop:current-mass-representation}
	设 \(U\subset\mathbb R^n\) 开，\(T\in\mathcal D_k(U)\) 具有局部有限质量。则存在唯一的 \(\Lambda_k\mathbb R^n\) 值 Radon 测度，代表 \(T\)。它可写成
	\begin{equation}\label{g3:eq:current-polar-representation}
		T(\omega)=\int_U\langle\omega(x),\vec T(x)\rangle\,\mathrm d\|T\|(x),
		\qquad |\vec T(x)|_{\mathrm m}=1\quad\|T\|\text{-几乎处处},
	\end{equation}
	这里 \(|\cdot|_{\mathrm m}\) 是形式余质量范数的对偶范数。其中正 Radon 测度 \(\|T\|\) 满足 \(\|T\|(V)=\mathbf M_V(T)\)，且 \(\operatorname{spt}T=\operatorname{spt}\|T\|\)。对 Borel 集 \(A\subset U\)，可定义限制流 \(T\llcorner A\)，其配对是在 \(A\) 上的积分，并有 \(\|T\llcorner A\|=\|T\|\llcorner A\)。
\end{proposition}
\symidx{\(\lVert T\rVert\)：局部有限质量流的质量测度}
\symidx{\(T\llcorner A\)：流 \(T\) 对 Borel 集 \(A\) 的限制}
\begin{proof}
	对每个坐标形式 \(\varphi\,\mathrm dx^I\) 应用\cref{g3:prop:bv-dual-variation}证明中使用的连续函数泛函测度表示：局部质量界使 \(\varphi\mapsto T(\varphi\,\mathrm dx^I)\) 关于一致范数连续，因而由带符号 Radon 测度 \(\mu_I\) 表示。相对紧开集上的表示由测试函数唯一性相容，得到全域的向量测度 \(\mu=\sum_I e_I\mu_I\)。

	令 \(\lambda=\sum_I|\mu_I|\)，写 \(\mu=h\lambda\)。以质量范数计算向量测度的全变差，得到 \(\|T\|=|h|_{\mathrm m}\lambda\)，并在 \(|h|_{\mathrm m}>0\) 处取 \(\vec T=h/|h|_{\mathrm m}\)。这一等式的证明与\cref{g3:lem:vector-measure-polar}相同：先对简单向量函数使用有限分割，再作 \(L^1(\lambda)\) 逼近。对偶上界
	\[
		|T(\omega)|\le\int\|\omega(x)\|_*\,\mathrm d\|T\|(x)
	\]
	直接成立。反向在有限维余质量单位球中取一个可数稠密族，逐点选择使配对逼近 \(|h|_{\mathrm m}\) 的元素；截成有限个取值，再由 Radon 正则性以紧支连续形式、继而以光滑形式逼近。投影到闭凸单位球或先将范数误差缩回一以内，保持测试界。故每个开集的对偶上确界正是 \(\|T\|\) 的质量。支集及限制的结论随之得到。
\end{proof}

这里 \(\vec T\) 只是一单位质量范数的向量场，还没有切平面或简单性保证。几何切向结构将在可整流流的定义及闭包证明中出现。对于 \(\partial_i\delta_a\)，取在 \(a\) 附近等于 \(\sin(j(x_i-a_i))\) 的单位模长紧支函数，配对绝对值为 \(j\)；因此一般流甚至不具有局部有限质量。

\begin{proposition}\label{g3:prop:current-mass-lsc}
	设 \(U\subset\mathbb R^n\) 开，\(0\le k\le n\)，\(T_j,T\in\mathcal D_k(U)\)。若 \(T_j\rightharpoonup T\)，即对每个紧支光滑 \(k\) 形式的配对收敛，则对每个开集 \(V\subset U\)，
	\[
		\mathbf M_V(T)\le\liminf_j\mathbf M_V(T_j).
	\]
	若各流局部质量有限，则对任意非负紧支连续函数 \(\eta\)，也有
	\[
		\int\eta\,\mathrm d\|T\|\le\liminf_j\int\eta\,\mathrm d\|T_j\|.
	\]
\end{proposition}
\begin{proof}
	对任意支于 \(V\) 且余质量至多为一的测试形式 \(\omega\)，都有
	\[
		T(\omega)=\lim_jT_j(\omega)\le\liminf_j\mathbf M_V(T_j).
	\]
	对 \(\omega\) 取上确界即可。加权结论将 \(\eta\) 写成开集 \(\{\eta>t\}\) 的层蛋糕积分，并使用 Fatou 引理。这个证明还表明：如果流列在每个相对紧开集上的质量一致有界，其弱极限自动具有局部有限质量。
\end{proof}

\begin{definition}\label{g3:def:normal-current}
	设 \(U\subset\mathbb R^n\) 开，\(0\le k\le n\)，\(T\in\mathcal D_k(U)\)，\(\partial T\) 为其边界，\(\mathbf M\) 为流的质量。若 \(T\) 与 \(\partial T\) 均有局部有限质量，称 \(T\) 为局部\term[zhenggui liu]{正规流}[normal current]；若
	\[
		\mathbf N(T):=\mathbf M(T)+\mathbf M(\partial T)<\infty,
	\]
	称 \(T\) 为正规流，记 \(T\in\mathcal N_k(U)\)。
\end{definition}
\symidx{\(\mathcal N_k(U)\)：流及其边界均有有限质量的正规流空间}

有限质量只控制流与形式的值配对，边界质量还控制它对外微分的作用。后一项正是下述连续变形估计的来源。对光滑函数 \(a\)，把乘法记作 \(T\llcorner a\)。若收缩约定为
\[
	(T\llcorner\mathrm da)(\eta)=T(\mathrm da\wedge\eta),
\]
则直接展开 \(\mathrm d(a\eta)\) 得到
\begin{equation}\label{g3:eq:current-product-rule}
	\partial(T\llcorner a)=(\partial T)\llcorner a-T\llcorner\mathrm da.
\end{equation}
光滑乘法因而保持局部正规性；任意 Borel 限制则未必保持边界质量。

\subsection{推前、乘积和同伦公式}
\label{g3:subsec:2201:03}
设 \(f:U\to V\) 光滑，且在 \(\operatorname{spt}T\) 上固有，即目标中每个紧集的原像与支集之交是 \(U\) 中的紧集。定义\term[liu de tuiqian]{流的推前}[pushforward of a current]
\[
	f_\#T(\omega)=T(\zeta f^*\omega),
\]
其中 \(\zeta\in C_c^\infty(U)\) 在 \(\operatorname{spt}T\cap f^{-1}(\operatorname{supp}\omega)\) 的邻域内为一。两个这样的截断之差在流的支集附近消失，故定义不依赖截断。紧支流的推前不需要另加固有性条件。由拉回与外微分相容，
\begin{equation}\label{g3:eq:current-pushforward-boundary}
	\partial(f_\#T)=f_\#(\partial T),\qquad
	(g\circ f)_\#T=g_\#(f_\#T).
\end{equation}
\symidx{\(f_\#T\)：流 \(T\) 沿映射 \(f\) 的推前}
若 \(T\) 有有限质量且 \(\|\mathrm Df\|\le L\) 于支集邻域，则
\begin{equation}\label{g3:eq:current-pushforward-mass}
	\mathbf M(f_\#T)\le L^k\mathbf M(T).
\end{equation}
这是因为 \(\|f^*\omega\|_*\le L^k\|\omega\|_*\)。这里一般只有不等式，多个原像的取向可能相消。

乘积流沿用“第一因子的方向在前”的约定。若 \(T\) 为 \(r\) 维流，\(S\) 为 \(s\) 维流，则对测试形式的坐标分量定义
\[
	(T\times S)(a(x,y)\,\mathrm dx^I\wedge\mathrm dy^J)
	=T\bigl(x\mapsto S(a(x,\cdot)\,\mathrm dy^J)\,\mathrm dx^I\bigr)
\]
当 \(|I|=r,|J|=s\)；其他双次数的分量取零。有限质量时，这就是两个向量测度的乘积与外积配对，因而可由 Fubini 定理计算。一般流的公式也成立：支集局部截断后，内层配对对外参数光滑，逐阶求导可由测试形式拓扑传入配对。

\begin{proposition}\label{g3:prop:current-homotopy}
	设 \(T\) 为 \(U\subset\mathbb R^n\) 中的紧支 \(k\) 维流，\(H:[0,1]\times U\to\mathbb R^q\) 光滑，记 \(H_t(x)=H(t,x)\)。则
	\begin{equation}\label{g3:eq:current-homotopy}
		H_{1\#}T-H_{0\#}T
		=\partial H_\#(\lBrack0,1\rBrack\times T)
		+H_\#(\lBrack0,1\rBrack\times\partial T).
	\end{equation}
	若 \(T\) 正规，且在相关支集上 \(|\partial_tH|\le a\)、\(\|\mathrm D_xH\|\le L\)，则
	\[
		\mathbf M\bigl(H_\#(\lBrack0,1\rBrack\times T)\bigr)
		\le aL^k\mathbf M(T).
	\]
	对 \(k\ge1\)，将 \(T\) 换为 \(\partial T\) 的相应上界为 \(aL^{k-1}\mathbf M(\partial T)\)；零维时该项为零。
\end{proposition}
\begin{proof}
	先在分量上展开外微分。\(x\) 方向的导数给出 \(\partial T\times S\)，\(y\) 方向的导数越过前面的 \(r\) 个微分，给出符号 \((-1)^r\)，所以
	\[
		\partial(T\times S)=\partial T\times S+(-1)^rT\times\partial S.
	\]
	对第一因子取区间，有
	\[
		\partial(\lBrack0,1\rBrack\times T)
		=\delta_1\times T-\delta_0\times T
		-\lBrack0,1\rBrack\times\partial T.
	\]
	推前并用式\eqref{g3:eq:current-pushforward-boundary}即得同伦公式。质量估计先在简单单位向量上计算：\(\partial_t\wedge\xi\) 的像是
	\(\partial_tH\wedge(\Lambda^k\mathrm D_xH)\xi\)，长度至多 \(aL^k\)。一般向量的质量范数可写为简单向量分解的长度和的下确界；这由余质量单位球的极集刻画得到，亦可用有限维分离定理验证。因此同一估计适用于一般向量测度。最后对 \(t\) 和 \(\|T\|\) 积分。
\end{proof}

对两个光滑 \(L\)-Lipschitz 映射 \(f,g\)，采用直线同伦。若在支集邻域内 \(|f-g|\le\delta\)，则对测试形式 \(\omega\)，
\begin{equation}\label{g3:eq:current-uniform-map-estimate}
\begin{aligned}
	|(f_\#T-g_\#T)(\omega)|
	&\le\delta L^k\mathbf M(T)\|\mathrm d\omega\|_{*,\infty}\\
	&\quad+\delta L^{k-1}\mathbf M(\partial T)\|\omega\|_{*,\infty}.
\end{aligned}
\end{equation}
这里及以下若 \(k=0\)，删除含 \(\partial T\) 的第二项。只控制映射的一致误差并不足以控制导数；正规流的边界质量通过同伦公式弥补了这一点。

\begin{theorem}\label{g3:thm:normal-lipschitz-pushforward}
	设 \(T\in\mathcal N_k(\mathbb R^n)\) 紧支，\(f:\mathbb R^n\to\mathbb R^q\) 为 Lipschitz 映射。则光滑卷积逼近 \(f_j\) 定义唯一的极限 \(f_\#T=\lim_jf_{j\#}T\)。任意在支集邻域一致趋于 \(f\)、Lipschitz 常数一致有界的光滑逼近给出同一流。推前与边界、复合相容，支集包含于 \(f(\operatorname{spt}T)\)，且
	\[
		\mathbf M(f_\#T)\le\operatorname{Lip}(f)^k\mathbf M(T).
	\]
	若 \(H:[0,1]\times\mathbb R^n\to\mathbb R^q\) 为 Lipschitz 同伦，则也满足式\eqref{g3:eq:current-homotopy}；若几乎处处 \(|\partial_tH|\le a\)、\(\|\mathrm D_xH\|\le L\)，对应同伦项的质量至多为 \(aL^k\mathbf M(T)\)。
\end{theorem}
\begin{proof}
	卷积 \(f_j=f*\rho_{\varepsilon_j}\) 保持 Lipschitz 常数，并在紧集上一致收敛。式\eqref{g3:eq:current-uniform-map-estimate}说明每个测试配对为 Cauchy 列。式\eqref{g3:eq:current-pushforward-mass}又使极限关于测试形式的一致范数有界，故定义一个有限质量流。对 \(\partial T\) 同样操作；因 \(\partial^2T=0\)，其配对同样收敛。对光滑推前的边界等式取极限，得到边界相容及正规性。支集结论由目标中与 \(f(\operatorname{spt}T)\) 不交的紧支测试形式检验得到。

	质量下半连续给出精确的 \(\operatorname{Lip}(f)^k\) 界。对于另一列光滑一致逼近，将它与卷积逼近代入同伦估计即可证明极限相同。复合相容可先固定一个光滑外映射，在内映射逼近中取极限，再逼近外映射；所有中间流的质量及边界质量由上述估计一致控制，故误差估计允许交换两个极限。同理，在 \([0,1]\times\operatorname{spt}T\) 的邻域内逼近 \(H\)，对区间乘积流应用刚建立的推前与边界相容，便把光滑同伦公式传到 Lipschitz 同伦。局部定义可先对支集附近作 Lipschitz 延拓；不同延拓在该邻域的推前一致。
\end{proof}

\begin{corollary}\label{g3:cor:normal-pushforward-finite}
	设 \(f:\mathbb R^n\to\mathbb R^q\) 为 Lipschitz 映射，\(0\le k\le n\)。
	\cref{g3:thm:normal-lipschitz-pushforward}的推前及边界相容可推广到任意 \(T\in\mathcal N_k(\mathbb R^n)\)，不要求紧支；支集包含于 \(\overline{f(\operatorname{spt}T)}\)。若 \(T\) 仅局部正规，而 \(f\) 在其支集上固有，则推前在目标中局部定义，保持局部正规性。
\end{corollary}
\begin{proof}
	取 \(0\le\chi_R\le1\)，在 \(B_R\) 上为一、支于 \(B_{2R}\)，且 \(|\nabla\chi_R|\le C/R\)。乘积公式说明
	\[
		\mathbf M(T\llcorner\chi_R-T)\to0,
		\mathbf M\bigl(\partial(T\llcorner\chi_R)-\partial T\bigr)\to0.
	\]
	对每个截断流使用紧支推前定理。质量估计应用于两个截断之差，使所得推前及其边界均为质量范数 Cauchy 列；其极限与截断选择无关，且继承质量界及边界公式。对局部正规且固有的情形，每个目标测试形式的相关原像只涉及源中的一个紧集，取在该集邻域为一的截断即可，交叠处由局部唯一性一致。
\end{proof}

还会用到局部版本的质量估计。设紧支正规流 \(T\) 的支集邻域中，映射的局部 Lipschitz 常数不超过一个连续函数 \(l(x)\)，则
\begin{equation}\label{g3:eq:current-weighted-pushforward}
	\mathbf M(f_\#T)\le\int l(x)^k\,\mathrm d\|T\|(x).
\end{equation}
若同伦在空间变量中的局部 Lipschitz 常数统一受 \(l(x)\) 控制，且时间方向速度受连续函数 \(a(x)\) 控制，则
\begin{equation}\label{g3:eq:current-weighted-homotopy}
	\mathbf M\bigl(H_\#(\lBrack0,1\rBrack\times T)\bigr)
	\le\int a(x)l(x)^k\,\mathrm d\|T\|(x).
\end{equation}
为证明这两式，对光滑映射先在每一点用外积估计，再对质量测度积分。对 Lipschitz 映射作卷积，小球内的导数界由该球上的 \(\sup l\) 控制；同伦的两个方向分别由 \(\sup a\)、\(\sup l\) 控制。在固定紧支邻域中这些函数有界，连续性及控制收敛使光滑估计传到上述积分界。后文的径向投影先删去奇点邻域，再应用这些连续局部控制。

\begin{proposition}\label{g3:prop:current-regularization-homotopy}
	设 \(0\le k\le n\)，\(T\) 为 \(\mathbb R^n\) 中紧支正规 \(k\) 维流，\(\rho_\varepsilon\) 为支于 \(B_\varepsilon\) 的非负单位积分光滑核。存在光滑系数流 \(\sigma_\varepsilon T\) 与正规 \((k+1)\) 维流 \(A_\varepsilon T\)，满足
	\[
		\sigma_\varepsilon T-T=\partial A_\varepsilon T+A_\varepsilon\partial T,
		\quad \mathbf M(\sigma_\varepsilon T)\le\mathbf M(T),
		\quad \mathbf M(A_\varepsilon T)\le\varepsilon\mathbf M(T).
	\]
	两者支集均包含于 \(\operatorname{spt}T+\overline B_\varepsilon\)，且 \(\partial\sigma_\varepsilon T=\sigma_\varepsilon\partial T\)。
\end{proposition}
\begin{proof}
	令 \(\tau_z(x)=x+z\)、\(H_z(t,x)=x+tz\)。对测试形式逐项定义
	\[
		\sigma_\varepsilon T=\int\rho_\varepsilon(z)\tau_{z\#}T\,\mathrm dz,
		\qquad A_\varepsilon T=\int\rho_\varepsilon(z)
		H_{z\#}(\lBrack0,1\rBrack\times T)\,\mathrm dz.
	\]
	前者的系数是向量测度与光滑核的卷积。后者由同伦质量界定义有限质量流。把同伦公式对 \(z\) 积分得到所述恒等式，并得到质量及支集估计。该恒等式使 \(\partial A_\varepsilon T\) 有有限质量，故 \(A_\varepsilon T\) 正规；边界相容由卷积与外微分相容得到。
\end{proof}

\subsection{光滑定向子流形及抵消现象}
\label{g3:subsec:2201:04}
设 \(M\subset\mathbb R^n\) 为 \(k\) 维光滑定向嵌入子流形，体积局部有限，取向由单位简单切向量 \(\tau_M\) 表示。它诱导积分流
\[
	\lBrack M\rBrack(\omega)
	=\int_M\langle\omega,\tau_M\rangle\,\mathrm d\mathcal H^k.
\]
局部面积公式保证这与坐标中的微分形式积分一致。其质量测度是 \(\mathcal H^k\llcorner M\)。若 \(M\) 紧且带光滑边界，由光滑 Stokes 公式，
\[
	\partial\lBrack M\rBrack=\lBrack\partial M\rBrack.
\]
这里边界采用外法向优先的取向；例如标准定向圆盘的边界圆周为逆时针。Stokes 的局部坐标及单位分解证明见第一卷第14.4节，当前使用的是紧光滑带边界情形。

\begin{example}\label{g3:exm:current-parallel-cancellation}
	在 \(\mathbb R^2\) 中令 \(Q_\varepsilon=(0,1)\times(0,\varepsilon)\)，取标准取向。边界流 \(T_\varepsilon=\partial\lBrack Q_\varepsilon\rBrack\) 的质量为 \(2+2\varepsilon\)。但对每个紧支光滑一形式，
	\[
		|T_\varepsilon(\omega)|
		=|\lBrack Q_\varepsilon\rBrack(\mathrm d\omega)|
		\le\varepsilon\|\mathrm d\omega\|_{*,\infty}\longrightarrow0.
	\]
	两条长边距离趋零、取向相反，最终抵消；两条短边的质量趋零。极限流为零，其质量严格小于质量极限。若只记录各条边上的无向长度，则两条长边会叠成重数二的线段。两种极限回答的几何问题不同。
\end{example}

对于同一光滑流形，反转取向使流乘以 \(-1\)，取整数份数使流乘以该整数。下一节把切平面、取向和重数推广到可数可整流集合，并以边界的相同结构选出整数流。
% 纲要条目排为编号小节标题，并进入目录与书签；说明段仍保留于纲要注释中。

% !TeX root = ../../Book3.tex
% 纲要标题：### 22.2 可整流流与整数流
% 纲要位置：几何学计划纲要.md，第 3018 行。
% 正文状态：初稿；尚未正式审核。

% 纲要内容开始
% 
% * 切平面、取向与整数重数
% * 边界质量和正规流
% * 整数流的定义
% * 取向与非定向模型的区别
% 
% 纲要内容结束

\section{可整流流与整数流}
\label{g3:sec:2202}

质量有限的流可以由弥散在整个空间中的向量测度表示，因而未必像一张曲面。要得到有向几何对象，还须规定质量集中于可整流集、向量沿切平面、重数属于指定的系数类。另一方面，一个有几何切向结构的流仍可能有无限边界质量。本节把这两种条件分开，再将它们结合为整数流。可整流表示的基础是\chapref{g3:ch:20}的 Lipschitz 分片与面积公式。

\subsection{切平面、取向与整数重数}
\label{g3:subsec:2202:01}
在一个可数 \(k\)-可整流 Borel 集上，可以先分成测度有限的 Lipschitz 片，在各片的几乎处处点取切平面。两个片的交叠处，切平面几乎处处相同：在参数域取交叠部分的密度一点，其微分决定相同的切向像空间。这使按片确定的切平面不依赖分片。当整个集合局部 \(\mathcal H^k\) 有限时，它正是第20章的近似切平面。

\begin{definition}\label{g3:def:rectifiable-current}
	设 \(U\subset\mathbb R^n\) 开，\(1\le k\le n\)。若 \(k\) 维流具有表示
	\begin{equation}\label{g3:eq:rectifiable-current-representation}
		T(\omega)=\int_M\theta(x)\langle\omega(x),\tau(x)\rangle\,
		\mathrm d\mathcal H^k(x),
	\end{equation}
	其中 \(M\subset U\) 为可数 \(k\)-可整流 Borel 集，\(\tau(x)\) 为其切平面的可测单位简单取向向量，\(\theta\) 为实值可测函数，且 \(\int_{M\cap K}|\theta|\,\mathrm d\mathcal H^k<\infty\) 对每个 \(K\Subset U\) 成立，则称 \(T\) 为实系数局部\term[kezhengliu liu]{可整流流}[rectifiable current]，记 \(T=\lBrack M,\tau,\theta\rBrack\)。若 \(\theta\) 几乎处处取整数值，称为整数重数可整流流。其有限质量全域类记为 \(\mathcal R_k(U)\)，局部类记为 \(\mathcal R_{k,\mathrm{loc}}(U)\)；本书的符号 \(\mathcal R_k\) 始终指整数重数类。
\end{definition}
\symidx{\(\mathcal R_k(U)\)：有限质量的整数重数可整流流类}
\symidx{\(\lBrack M,\tau,\theta\rBrack\)：可整流流的位置、取向与重数表示}

可以删去 \(\theta=0\) 的点，把负号吸收到 \(\tau\) 中，从而令 \(\theta>0\)。整数重数时 \(\theta\ge1\)，所以局部有限质量自动使有效承载集局部 \(\mathcal H^k\) 有限。取向反转与重数变号给出同一流；只有它们的乘积有确定意义。由质量测度表示，
\begin{equation}\label{g3:eq:rectifiable-current-mass}
	\|T\|=|\theta|\,\mathcal H^k\llcorner M.
\end{equation}

零维时，定义实系数可整流流为可数点流之和 \(\sum_i a_i\delta_{x_i}\)，要求每个紧集内 \(\sum|a_i|<\infty\)。整数重数时 \(a_i\in\mathbb Z\)，合并相同位置并删去零系数后，紧集中只有有限个点；若总质量有限，整个和式只有有限项。一般零维有限质量流却可以是非原子的 Radon 测度。

\begin{proposition}\label{g3:prop:rectifiable-lipschitz-pushforward}
	设 \(T=\lBrack M,\tau,\theta\rBrack\) 为 \(\mathbb R^n\) 中有限质量的实系数可整流 \(k\) 维流，\(f:\mathbb R^n\to\mathbb R^q\) 为 Lipschitz 映射。记 \(\mathrm D^Mf\) 为沿近似切平面的微分，\(J_k^Mf\) 为其 \(k\) 维 Jacobian 面积因子。则
	\begin{equation}\label{g3:eq:rectifiable-pushforward-formula}
		f_\#T(\omega)=\int_M\theta(x)
		\big\langle\omega(f(x)),(\Lambda^k\mathrm D^Mf(x))\tau(x)\big\rangle
		\mathrm d\mathcal H^k(x)
	\end{equation}
	定义一个实系数可整流流。若 \(T\) 为整数重数，则 \(f_\#T\) 也为整数重数。质量满足
	\[
		\mathbf M(f_\#T)\le\int_M|\theta|J_k^Mf\,\mathrm d\mathcal H^k
		\le\operatorname{Lip}(f)^k\mathbf M(T).
	\]
	若 \(T\) 还正规，上式与\cref{g3:thm:normal-lipschitz-pushforward,g3:cor:normal-pushforward-finite}的定义一致。
\end{proposition}
\begin{proof}
	可整流集上的微分 \(\mathrm D^Mf\) 是各 Lipschitz 参数片上的微分，经参数的单射微分识别到切平面所得；交叠处由密度一点的链式法则一致。将满秩部分分成可数双 Lipschitz 片，秩不足部分的像由面积公式具有零 \(\mathcal H^k\) 测度。像的各片也可整流，在交叠处取一个共同的可测单位取向 \(\tau'\)。对满秩原像，
	\[
		(\Lambda^k\mathrm D^Mf(x))\tau(x)
		=\varepsilon(x)J_k^Mf(x)\tau'(f(x)),\qquad \varepsilon(x)\in\{-1,1\}.
	\]
	第20章带权面积公式将式\eqref{g3:eq:rectifiable-pushforward-formula}改写为像上的积分，重数为
	\[
		\theta'(y)=\sum_{x\in M\cap f^{-1}(y),\,J_k^Mf(x)>0}
		\varepsilon(x)\theta(x).
	\]
	该级数几乎处处绝对收敛，因为绝对值之和的积分不超过 \(\operatorname{Lip}(f)^k\int_M|\theta|\)。整数输入时，非零整数项的绝对可和迫使几乎每根纤维只有有限个非零贡献，所以和仍为整数。三角不等式给出质量估计。

	还须比较这个切向公式与光滑逼近定义。所需解析事实是：若 \(D\subset\mathbb R^k\) 开，\(g_j:D\to\mathbb R^q\) 一致 Lipschitz，并在紧集上一致趋于 \(g\)，则每个 \(k\) 阶微分子式在局部 \(L^\infty\) 的弱星意义下趋于相应极限。证明对阶数归纳。对一个紧支光滑函数 \(\eta\)，积分分部给出
	\[
		\int\eta\,\mathrm dg_j^{i_1}\wedge\cdots\wedge\mathrm dg_j^{i_\ell}
		\wedge\mathrm dx^J
		=-\int g_j^{i_1}\,\mathrm d\eta\wedge\mathrm dg_j^{i_2}
		\wedge\cdots\wedge\mathrm dg_j^{i_\ell}\wedge\mathrm dx^J,
	\]
	其中 \(J\) 补足总次数。\(\ell=1\) 是弱导数定义；一般情形用 \(g_j^{i_1}\) 的一致收敛及归纳假设，先把连续系数一致逼近为光滑系数，得到配对收敛。一致的子式界再把测试系数从光滑函数扩大到 \(L^1\)。Lipschitz 函数的积分分部可先卷积后取极限，按相同归纳成立。

	现在取 \(M\) 的一个双 Lipschitz 参数片 \(\psi:A\subset\mathbb R^k\to M\)，将 \(\psi\) 按坐标延拓到 \(\mathbb R^k\)，再取 \(g_j=f_j\circ\psi\)。双 Lipschitz 性给出参数 Jacobian 的正下界，所以面积换元后 \(\theta\circ\psi\) 在 \(A\) 上可积。当前配对化成参数域中带 \(L^1\) 系数的子式积分；系数 \(\omega(f_j\circ\psi)\) 一致趋于极限，\(\mathbf1_A\) 及可积重数保留在 \(L^1\) 系数中。故每个片上的积分趋于切向公式。先保留有限个有界片，使其余片的 \(|\theta|\) 积分任意小，再利用统一 Lipschitz 质量界处理尾项，即得整体收敛。正规紧支情形由唯一性得到一致；一般有限质量正规流再作\cref{g3:cor:normal-pushforward-finite}中的截断，利用质量范数收敛得到相同结论。
\end{proof}

特别地，Lipschitz 图嵌入 \(F(x)=(x,f(x))\) 将可整流流送到其图上，取向是 \((\Lambda^k\mathrm D^MF)\tau\) 的归一化。因第一坐标保留 \(x\)，不存在不同原像抵消。这个构造使下一节能够把一般 Lipschitz 函数的切片化为图上的光滑高度切片。

\begin{lemma}\label{g3:lem:integer-rectifiable-mass-closure}
	设 \(T_j\in\mathcal R_k(\mathbb R^n)\)，\(T\) 为有限质量流，且 \(\mathbf M(T_j-T)\to0\)。则 \(T\in\mathcal R_k(\mathbb R^n)\)。
\end{lemma}
\begin{proof}
	设 \(M_j\) 承载 \(T_j\)，则 \(M=\bigcup_jM_j\) 为可数可整流集，且质量范数收敛说明 \(\|T\|(\mathbb R^n\setminus M)=0\)。把 \(M\) 分成互不交的有限测度参数片；片间交叠的切平面几乎处处相容，故可在并集上选择一个共同的可测取向 \(\tau\)。每个 \(T_j\) 在此取向下的系数 \(\theta_j\) 为整数，片外补零，并且
	\[
		\mathbf M(T_i-T_j)=\int_M|\theta_i-\theta_j|\,\mathrm d\mathcal H^k.
	\]
	因此 \(\theta_j\) 是 \(L^1(M,\mathcal H^k)\) 中的 Cauchy 列。其极限 \(\theta\) 沿子列几乎处处为整数，且 \(\int|\theta|<\infty\)。测试配对确定 \(T=\lBrack M,\tau,\theta\rBrack\)。删去 \(\theta=0\) 的部分后，有效承载集的测度不超过 \(\int|\theta|\)。零维时同样是可数点集上的整数 \(\ell^1\) 极限。
\end{proof}

\subsection{边界质量和正规流}
\label{g3:subsec:2202:02}
\cref{g3:def:normal-current}要求流及其边界的质量都有限。这个条件并不规定重数为整数，也不规定质量集中于与流同维的集合。例如取 \(k<n\)、非零 \(a\in C_c^\infty(\mathbb R^n)\)，令
\[
	T(\omega)=\int_{\mathbb R^n}a(x)\langle\omega(x),e_1\wedge\cdots\wedge e_k\rangle\,\mathrm dx.
\]
对各坐标积分分部，\(\partial T\) 的系数是 \(a\) 的一阶导数，故 \(T\) 正规。但其质量对 \(n\) 维 Lebesgue 测度绝对连续，任何可数 \(k\)-可整流集均为 \(n\) 维零测集，所以这个非零流不是 \(k\) 维可整流流。

正规性仍强制一些切向约束。下面两条结果将用于立方网格的面上。

\begin{lemma}\label{g3:lem:normal-current-no-lower-skeleton}
	设 \(T\) 为开集 \(U\subset\mathbb R^n\) 中局部正规的 \(k\) 维流，\(k\ge1\)。对任意仿射超平面 \(H\)，限制向量测度 \(T\llcorner H\) 只含沿 \(H\) 的 \(k\) 向量。特别地，若仿射平面 \(P\) 的维数小于 \(k\)，则 \(\|T\|(P\cap U)=0\)；有限个这样的平面之并也不承载质量。
\end{lemma}
\begin{proof}
	旋转、平移后令 \(H=\{x_j=0\}\)。取 \(\rho\in C_c^\infty(\mathbb R)\)，在零附近为一，令 \(a_\varepsilon(s)=s\rho(s/\varepsilon)\)。则 \(\|a_\varepsilon\|_\infty\le C\varepsilon\)，导数一致有界，且逐点趋于 \(\mathbf1_{\{0\}}\)。若多重指标 \(I\) 含 \(j\)，以
	\(a_\varepsilon(x_j)\eta(x)\,\mathrm dx^{I\setminus j}\)
	测试边界。除去含 \(a'_\varepsilon\) 的一项，其余两项分别由 \(C\varepsilon\|\partial T\|(\operatorname{supp}\eta)\) 和 \(C\varepsilon\|T\|(\operatorname{supp}\eta)\) 控制。由向量测度表示与控制收敛，令 \(\varepsilon\downarrow0\)，得到 \(\int_H\eta\,\mathrm d\mu_I=0\)。因此所有含法向指标的系数测度在 \(H\) 上消失。

	若 \(P\) 的维数为 \(d<k\)，取适配坐标使 \(P\subset\{x_{d+1}=\cdots=x_n=0\}\)。每个 \(k\) 重指标都含某个法向指标，刚才的结论使相应测度在 \(P\) 上为零。所有分量均为零，因而质量也为零。
\end{proof}

\begin{proposition}\label{g3:prop:normal-plane-bv}
	设 \(P\subset\mathbb R^n\) 为定向 \(k\) 维仿射平面，\(1\le k\le n\)，\(W\subset\mathbb R^n\) 开。若局部正规流 \(T\in\mathcal D_k(W)\) 支于 \(P\cap W\)，则存在唯一 \(u\in BV_{\mathrm{loc}}(P\cap W)\)，使
	\[
		T=u\lBrack P\rBrack,
		\qquad \|T\|=|u|\mathcal H^k\llcorner(P\cap W),
		\qquad \|\partial T\|=|D_Pu|.
	\]
	最后一个等式以 \(P\) 的等距坐标将 BV 变差测度视为环境中的测度。\(k=0\) 时，支于单点的有限质量流是该点流的实倍数。
\end{proposition}
\begin{proof}
	先取适配坐标 \(P=\mathbb R^k\times\{0\}\)。上一引理排除了所有带法向指标的系数，所以 \(T=e_1\wedge\cdots\wedge e_k\,\mu\)，其中 \(\mu\) 是 \(P\) 上的带符号局部有限测度。对切向测试形式计算边界，说明每个分布偏导 \(D_i\mu\) 都是局部有限测度；它们正是 \(\partial T\) 的各切向分量，符号由删去第 \(i\) 个基向量决定。

	说明 \(\mu\) 实际具有 BV 密度。在一个相对紧切向区域附近先乘紧支光滑截断，把 \(\mu\) 及其偏导变成有限测度。卷积 \(u_\varepsilon=\rho_\varepsilon*\mu\) 的 \(L^1\) 范数和梯度 \(L^1\) 范数由这些测度的质量一致控制。\cref{g3:thm:bv-local-compactness}给出子列局部 \(L^1\) 趋于 \(u\in BV_{\mathrm{loc}}\)。另一方面卷积在分布意义下趋于 \(\mu\)，所以 \(\mu=u\mathcal L^k\)。在交叠区域中由测度唯一性粘合。

	最后，把向量 \(D_Pu\) 通过定向的切向 Hodge 对偶识别为 \((k-1)\) 向量；这个识别保持长度，且在 \(k\) 维空间中的 \((k-1)\) 向量均简单。因此边界质量恰为 BV 全变差，流质量则为 \(|u|\mathcal H^k\)。唯一性由测试顶次形式立即得到。
\end{proof}

\begin{theorem}[常数定理]\label{g3:thm:current-constancy}
	设 \(V\) 为定向 \(k\) 维仿射平面中的连通开集，\(k\ge1\)。若 \(T\) 为其中的局部正规顶维流且 \(\partial T=0\)，则 \(T=c\lBrack V\rBrack\)，其中 \(c\in\mathbb R\) 为常数。若 \(T\) 为整数重数可整流流，则 \(c\in\mathbb Z\)。
\end{theorem}
\begin{proof}
	上一命题给出 \(T=u\lBrack V\rBrack\)，且 \(Du=0\)。在任意较小球内卷积，所得光滑函数梯度为零，故为常数。令卷积尺度趋零，得到 \(u\) 在该球中几乎处处为常数。连通开集中的任意两点可用有限条相交小球连接，所以这些常数相同。整数重数情形，\(u\) 几乎处处为整数，故常数为整数。
\end{proof}

\subsection{整数流的定义}
\label{g3:subsec:2202:03}
一张定向多面体的每个面带有整数份数，其边界仍由整数份数的低维面组成。一般可整流对象也需要边界具有这一性质。

\begin{definition}\label{g3:def:integer-current}
	设 \(U\subset\mathbb R^n\) 开，\(0\le k\le n\)，\(\mathcal R_j(U)\) 表示有限质量的整数重数可整流 \(j\) 维流类。若 \(T\in\mathcal R_k(U)\) 且 \(\partial T\in\mathcal R_{k-1}(U)\)，称 \(T\) 为\term[zhengshu liu]{整数流}[integral current]，记为 \(\mathcal I_k(U)\)。零维时取 \(\mathcal I_0(U)=\mathcal R_0(U)\)，因为其边界为零。将两处质量有限改为局部质量有限，得到局部整数流类 \(\mathcal I_{k,\mathrm{loc}}(U)\)。
\end{definition}
\symidx{\(\mathcal I_k(U)\)：有限质量的 \(k\) 维整数流类}

每个整数流都正规，而一般正规流的重数可以是实数，甚至没有可整流承载集。上述定义同时规定 \(T\) 与 \(\partial T\) 的整数可整流性；后面的闭包定理会分别恢复极限与极限边界的这两项结构。

\begin{proposition}\label{g3:prop:integer-current-operations}
	设 \(U\subset\mathbb R^n\) 开，\(0\le k\le n\)。下述流的整数性均包含自身及边界的整数重数可整流性和有限质量。
	整数流对加法、相反数、整数倍数、边界运算封闭。若 \(T\in\mathcal I_k(U)\) 紧支，\(f:U\to\mathbb R^q\) 为 Lipschitz 映射，则 \(f_\#T\) 为整数流。若 \(T\in\mathcal I_k(\mathbb R^n)\)、\(a<b\)，则 \(\lBrack a,b\rBrack\times T\) 为整数流。
\end{proposition}
\begin{proof}
	将有限个可整流承载集取并，交叠处采用共同取向，重数按带符号整数相加；质量由三角不等式控制。边界同样满足这些性质，且 \(\partial^2T=0\)，所以前四项成立。推前由\cref{g3:prop:rectifiable-lipschitz-pushforward}保持整数可整流表示，对边界再用式\eqref{g3:eq:current-pushforward-boundary}。

	区间乘积的承载集是 \((a,b)\times M\)。各参数片乘以区间仍为 Lipschitz 片，切向量为 \(e_t\wedge\tau\)，重数仍为 \(\theta\)，且质量为 \((b-a)\mathbf M(T)\)。乘积边界公式给出
	\[
		\partial(\lBrack a,b\rBrack\times T)
		=\delta_b\times T-\delta_a\times T-\lBrack a,b\rBrack\times\partial T.
	\]
	右侧均为有限质量整数可整流流，故乘积为整数流。
\end{proof}

\begin{example}\label{g3:exm:finite-perimeter-integral-current}
	设 \(E\subset\mathbb R^n\) 可测，\(|E|<\infty\)、\(P(E)<\infty\)。用标准取向定义 \(T_E=\lBrack E,e_1\wedge\cdots\wedge e_n,1\rBrack\)。令 \(\mathrm{vol}=\mathrm dx^1\wedge\cdots\wedge\mathrm dx^n\)，每个紧支 \((n-1)\) 形式唯一写成 \(\iota_X\mathrm{vol}\)，且
	\[
		\partial T_E(\iota_X\mathrm{vol})
		=\int_E\operatorname{div}X\,\mathrm dx
		=\int_{\partial^*E}X\cdot\nu_E\,\mathrm d\mathcal H^{n-1}.
	\]
	末式是\cref{g3:thm:finite-perimeter-gauss-green}。在 \(\partial^*E\) 上选单位切向量 \(\tau_{\partial E}\)，使
	\(\nu_E\wedge\tau_{\partial E}=e_1\wedge\cdots\wedge e_n\)。约化边界结构定理给出
	\[
		\partial T_E=\lBrack\partial^*E,\tau_{\partial E},1\rBrack,
		\qquad \mathbf M(T_E)=|E|,\qquad \mathbf M(\partial T_E)=P(E).
	\]
	因此 \(T_E\) 是整数流。向量测度 \(D\chi_E=-\nu_E|D\chi_E|\) 的负号与此处外法向优先的边界取向相容：二者分别通过 \(-\int_E\operatorname{div}X\) 和 \(\int_E\operatorname{div}X\) 定义。
\end{example}

在最高维，前节的 BV 表示给出更一般的字典：\(T=u\lBrack\mathbb R^n\rBrack\) 正规当且仅当 \(u\in BV(\mathbb R^n)\)，且
\[
	\mathbf M(T)=\|u\|_{L^1},\qquad \mathbf M(\partial T)=|Du|(\mathbb R^n).
\]
若 \(u\) 为整数值，层蛋糕分解把正负部分写成整数层集的和。具体地，令
\[
	E_j^+=\{u\ge j\},\qquad E_j^-=\{u\le-j\},\qquad j=1,2,\ldots.
\]
BV 余面积公式及 \(u\) 的整数值性质给出
\[
	u=\sum_{j\ge1}(\chi_{E_j^+}-\chi_{E_j^-}),\qquad
	\sum_{j\ge1}\bigl(P(E_j^+)+P(E_j^-)\bigr)=|Du|(\mathbb R^n).
\]
每个集合体积有限且周长有限。上述例子的边界流按质量范数可求和，\cref{g3:lem:integer-rectifiable-mass-closure}说明和的边界仍为整数可整流流。故整数值 BV 函数产生顶维整数流。这里 BV 余面积公式使用的是整段水平区间的积分，而不是假定每个指定实水平都为良水平。

反过来，整数可整流性本身不控制边界。例如在 \((0,1)\) 内取两两不交的区间 \((a_j,b_j)\)，总长度有限，所有端点不同且趋于同一个内点。流 \(T=\sum_j\lBrack a_j,b_j\rBrack\) 质量有限，是整数重数可整流流。对任意有限组端点，可选模长不超过一的光滑测试函数分别取对应的正负号，且避开其余端点。因此 \(\mathbf M(\partial T)\ge2N\) 对每个 \(N\) 成立，其边界质量无穷。

\subsection{取向与非定向模型的区别}
\label{g3:subsec:2202:04}
整数重数可整流流允许选取可测取向，不要求全局连续。这个自由并不消除定向障碍：取向跳变通常会在边界流中留下内部项。

\begin{proposition}\label{g3:prop:current-orientation-obstruction}
	设 \(M\subset\mathbb R^n\) 为连通光滑 \(k\) 维嵌入流形，可带边界。若一个局部正规流 \(T\) 在 \(M\) 上具有模长为一的整数重数表示，并且 \(\operatorname{spt}\partial T\subset\partial M\)，则 \(M\) 可定向。因而非定向光滑流形不能以单位整数重数表示成只有几何边界的流。
\end{proposition}
\begin{proof}
	在 \(M\) 的任意内点附近取光滑坐标，把一小片展平为 \(\mathbb R^k\) 中的连通开集；可将图坐标扩为环境中的局部光滑坐标变化。推前保持边界等式。因为此邻域没有边界流，\cref{g3:thm:current-constancy}说明表示密度为常数。单位重数条件使常数为 \(+1\) 或 \(-1\)。因此流在每个坐标片上确定一个连续取向。

	在两片交叠处，它们给出同一个非零流；若取向不相容，同一正体积小片上的测试配对将相反，矛盾。故局部取向相容，粘成整个内部的取向。带边界坐标的内侧为连通半球，这个取向唯一连续延伸到边界。因此 \(M\) 可定向。
\end{proof}

以 Möbius 带为例，把它沿一条横截线切开，得到一个可定向矩形。将矩形的积分流推回带上，矩形两条被识别的边因粘合反转了横向坐标，在流边界中产生同向的两份横截线。这个内部边界正记录了无法连续选取向的障碍。只用面积而不记录取向的模型，或把系数改成模二，会有不同的相消规则；本章固定使用实数配对中的整数重数。

同样，一个几何集合不自动确定流。重合的两张同向片产生重数二，反向片产生零；相同的支集也可以配不同整数重数。质量、取向和边界必须同时保留，才能在后面讨论固定边界的极小化问题。
% 纲要条目排为编号小节标题，并进入目录与书签；说明段仍保留于纲要注释中。

% !TeX root = ../../Book3.tex
% 纲要标题：### 22.3 切片、锥与平坦范数
% 纲要位置：几何学计划纲要.md，第 3025 行。
% 正文状态：初稿；尚未正式审核。

% 纲要内容开始
% 
% * Lipschitz 切片及边界公式
% * 锥构造和填充估计
% * 平坦范数与弱收敛
% * 质量范数收敛、平坦收敛与弱收敛
% 
% 纲要内容结束

\section{切片、锥与平坦范数}
\label{g3:sec:2203}

把一条有向曲线限制在半空间内，边界中会出现它穿过切割面的交点。高维流的切片记录同一种现象：它是截断后新出现的边界。切片把维数降低，锥把维数升高，而平坦范数同时允许一个小的本体误差与一个小填充的边界。这三种操作共同为后面的变形和紧性提供定量控制。

\subsection{Lipschitz 切片及边界公式}
\label{g3:subsec:2203:01}
先设 \(T\) 是 \(\mathbb R^n\) 中的正规 \(k\) 维流，\(k\ge1\)，\(f:\mathbb R^n\to\mathbb R\) 为 Lipschitz 函数。记
\[
	T_t=T\llcorner\{f<t\},\qquad
	S_t=\partial T_t-(\partial T)\llcorner\{f<t\}.
\]
右端首先是测试形式上的流。我们要证明：几乎每个 \(t\) 对应的 \(S_t\) 实际有有限质量，而且质量的参数积分由 \(T\) 控制。这样，截断所增加的边界便不会失控。上述约定与有限周长集合的外法向取向一致；切片的定义和以下公式可对照\cite[第28节，第156--162页]{Simon1983}。

\begin{theorem}\label{g3:thm:normal-current-slicing}
	设 \(T\in\mathcal N_k(\mathbb R^n)\)，\(1\le k\le n\)，\(f:\mathbb R^n\to\mathbb R\) 为 \(L\)-Lipschitz 函数。对几乎所有 \(t\)，\term[qiepian]{切片}[slice]
	\begin{equation}\label{g3:eq:current-slice-definition}
		\langle T,f,t\rangle
		=\partial(T\llcorner\{f<t\})-(\partial T)\llcorner\{f<t\}
	\end{equation}
\symidx{\(\langle T,f,t\rangle\)：流 \(T\) 沿实值 Lipschitz 函数水平集的切片}
	属于 \(\mathcal N_{k-1}(\mathbb R^n)\)，且有
	\[
		\operatorname{spt}\langle T,f,t\rangle
		\subset\operatorname{spt}T\cap f^{-1}(t),
		\qquad
		\partial\langle T,f,t\rangle=-\langle\partial T,f,t\rangle.
	\]
	这里约定零维流的切片及边界均为零。对任意 Borel 集 \(B\subset\mathbb R^n\) 及 \(-\infty\le a<b\le\infty\)，有
	\begin{equation}\label{g3:eq:normal-slice-mass}
		\int_a^b\|\langle T,f,t\rangle\|(B)\,\mathrm dt
		\le L\|T\|(B\cap\{a<f<b\}).
	\end{equation}
	特别地，\(\int_{\mathbb R}\mathbf N(\langle T,f,t\rangle)\,\mathrm dt\le L\mathbf N(T)\)，且几乎所有截断 \(T\llcorner\{f<t\}\) 仍为正规流。
\end{theorem}
\begin{proof}
	证明的要点是先对切片的全部测试值建立一个积分估计，再由可数稠密的测试族取得质量。这样不必预先假定切片可以由某个曲面表示。

	先假定 \(f\) 光滑。对 \(\omega\in\mathcal D^{k-1}(\mathbb R^n)\) 和 \(h\in C_c^\infty(\mathbb R)\)，令 \(H(s)=\int_s^\infty h(t)\,\mathrm dt\)。按质量测度作 Fubini 换序，并在 \(\operatorname{spt}\omega\) 的邻域乘一个等于一的截断，得到
	\begin{equation}\label{g3:eq:current-slice-pairing}
	\begin{aligned}
		\int_{\mathbb R}h(t)S_t(\omega)\,\mathrm dt
		&=T(H(f)\,\mathrm d\omega)-\partial T(H(f)\omega)\\
		&=T\bigl(h(f)\,\mathrm df\wedge\omega\bigr).
	\end{aligned}
	\end{equation}
	因 \(|\mathrm df\wedge\omega|_*\le L|\omega|_*\)，右端由质量控制。这里 \(|\cdot|_*\) 表示形式的余质量范数。

	固定开集 \(U\) 及有限区间 \((a,b)\)，在支于 \(U\)、余质量至多一的光滑 \((k-1)\) 形式中取可数的 \(C^1\) 稠密族 \(\omega_i\)，包含零形式和各形式的相反数。每个 \(t\mapsto S_t(\omega_i)\) 可测。在有限个测试值中按取得最大值的最小下标划分参数区间；用彼此不交的紧子集及不交支撑的光滑函数 \(h_i\) 逼近这些可测片，并保持 \(0\le\sum_i h_i\le1\)、\(\operatorname{spt}h_i\subset(a,b)\)。这种逼近可以在有限测度 \(\sum_i|S_t(\omega_i)|\,\mathrm dt\) 下完成。对式\eqref{g3:eq:current-slice-pairing}求和，便有
	\[
		\int_a^b\max_{i\le N}S_t(\omega_i)\,\mathrm dt
		\le L\|T\|(U\cap\{a<f<b\}).
	\]
	令 \(N\to\infty\)，单调收敛说明切片的质量上确界几乎处处有限，并给出式\eqref{g3:eq:normal-slice-mass}的开集版本。对可数个开集基使用这些测试族，再用测度的单调类论证，\(t\mapsto\|S_t\|(B)\) 对每个 Borel 集 \(B\) 可测；其积分定义的测度被 \(L\|T\|\) 控制。由正则性，估计推广到全部 Borel 集。

	对一般 Lipschitz \(f\)，令 \(f_j=f*\rho_{\varepsilon_j}\)。则 \(f_j\to f\) 一致，且 \(\operatorname{Lip}f_j\le L\)。有限测度 \(\|T\|+\|\partial T\|\) 只可能对可数个水平集赋予正质量。因此对几乎所有 \(t\)，截断的两个项都在测试形式上收敛，得到
	\[
		S_t^{(j)}\rightharpoonup
		\partial(T\llcorner\{f<t\})-(\partial T)\llcorner\{f<t\}=S_t.
	\]
	质量是测试值上确界，故在这个弱收敛下下半连续。对开集 \(U\) 使用 Fatou 引理，先在端点水平集的质量均为零时令 \(j\to\infty\)，得到相同的带状估计；一般 \(a,b\) 可由内部的非原子水平递增穷尽。再次以正则性推广到 Borel 集。可测性仍由截断定义和可数测试族得到。

	若测试形式支于 \(\{f<t\}\) 内，则截断公式中的两个项相消；若支于 \(\{f>t\}\)，则两项都为零。这给出切片的支集断言。对定义施加边界算子，并使用 \(\partial^2T=0\)，得到
	\[
		\partial S_t=-\partial\bigl((\partial T)\llcorner\{f<t\}\bigr)
		=-\langle\partial T,f,t\rangle.
	\]
	把已经证明的质量估计用于 \(\partial T\)，便知右端几乎处处有限质量，且有相应积分界；\(k=1\) 时两边本来就是零。这证明 \(S_t\) 正规及 \(\mathbf N\) 的估计。最后
	\(\partial T_t=(\partial T)\llcorner\{f<t\}+S_t\)，所以几乎所有 \(T_t\) 也正规。
\end{proof}

局部正规流也有同样结论：先在相对紧开集附近乘光滑截断，应用定理，再用可数穷尽删去共同的零测参数集。在两个截断均等于一的区域内，定义完全相同，故所得局部切片可以粘合。后面按距离函数选择切割半径时，只需式\eqref{g3:eq:normal-slice-mass}；个别异常半径不会影响选择。

对于可整流流，切片还可以直接读成水平集上的有向积分。若 \(M\) 可数 \(k\)-可整流，\(f\) 为 Lipschitz 函数，则 \(\mathrm d_Mf\) 表示其沿近似切平面的微分，\(\nabla_Mf\) 是对应切向量。由双 Lipschitz 参数片定义它时，在参数域的密度点求微分；第20章的延拓无关性保证它不依赖片的选择。我们采用收缩约定
\[
	\langle\tau\mathbin{\lrcorner}\alpha,\omega\rangle
	=\langle\tau,\alpha\wedge\omega\rangle.
\]
当 \(\tau\) 是单位简单 \(k\) 向量时，\(\tau\mathbin{\lrcorner}\mathrm d_Mf\) 的长度等于 \(|\nabla_Mf|\)。

\begin{theorem}\label{g3:thm:integer-current-slicing}
	设 \(T=\lBrack M,\tau,\theta\rBrack\) 是 \(\mathbb R^n\) 中有限质量的可整流 \(k\) 维流，\(k\ge1\)，\(\theta\) 为正的可测实重数，并设 \(\mathbf M(\partial T)<\infty\)。令 \(f:\mathbb R^n\to\mathbb R\) 为 Lipschitz 函数，\(M^+=\{x\in M:|\nabla_Mf(x)|>0\}\)。则对几乎所有 \(t\)，有
	\begin{equation}\label{g3:eq:rectifiable-slice-representation}
		\langle T,f,t\rangle
		=\left\lBrack M^+\cap f^{-1}(t),
		\frac{\tau\mathbin{\lrcorner}\mathrm d_Mf}{|\nabla_Mf|},\theta\right\rBrack.
	\end{equation}
	右侧为可整流 \((k-1)\) 维流；若 \(\theta\) 几乎处处为整数，则切片也整数可整流。若进一步 \(T\in\mathcal I_k(\mathbb R^n)\)，则几乎所有切片属于 \(\mathcal I_{k-1}\)，而 \(T\llcorner\{f<t\}\) 属于 \(\mathcal I_k\)。
\end{theorem}
\begin{proof}
	先把第20章的余面积公式移到 \(M\) 上。由\cref{g3:prop:rectifiable-bilipschitz-cover}，把 \(M\) 分为互不交的双 Lipschitz 参数片 \(\varphi:A\subset\mathbb R^k\to M\)。对 \(f\circ\varphi\) 使用\cref{g3:thm:lipschitz-coarea-formula}，并在参数片及其纤维上分别使用加权面积公式。在线性代数层面，\(\mathrm D\varphi\) 将 \(\ker\mathrm D(f\circ\varphi)\) 映到 \(T_xM\cap\ker\mathrm d_Mf\)，且
	\[
		J_k\varphi\,|\nabla_Mf|
		=J_{k-1}\bigl(\mathrm D\varphi|_{\ker\mathrm D(f\circ\varphi)}\bigr)
		\,|\nabla(f\circ\varphi)|.
	\]
	这可在核内取正交基，再补一个法向向量，以平行多面体的底面积乘高验证。于是对非负可测 \(g\)，逐片求和得到
	\begin{equation}\label{g3:eq:rectifiable-coarea-scalar}
		\int_M g|\nabla_Mf|\,\mathrm d\mathcal H^k
		=\int_{\mathbb R}\left(\int_{M\cap f^{-1}(t)}g\,\mathrm d\mathcal H^{k-1}\right)\mathrm dt.
	\end{equation}
	第20章对满秩参数片的纤维构造同时给出几乎所有水平集的可整流性，及其切平面 \(T_xM\cap\ker\mathrm d_Mf\)。秩亏部分和参数零测余项的纤维几乎处处为零测集。取 \(g=\theta\) 可知右侧有有限的重数积分。

	先令 \(f\) 光滑。把式\eqref{g3:eq:current-slice-pairing}中的右端写成 \(M\) 上积分，再使用式\eqref{g3:eq:rectifiable-coarea-scalar}，得到
	\[
	\begin{aligned}
		\int h(t)\langle T,f,t\rangle(\omega)\,\mathrm dt
		&=\int_M h(f)\theta\langle\tau,\mathrm d_Mf\wedge\omega\rangle\,\mathrm d\mathcal H^k\\
		&=\int h(t)\left(\int_{M^+\cap f^{-1}(t)}
		\theta\left\langle\frac{\tau\mathbin{\lrcorner}\mathrm d_Mf}{|\nabla_Mf|},\omega\right\rangle
		\,\mathrm d\mathcal H^{k-1}\right)\mathrm dt.
	\end{aligned}
	\]
	因此两个切片对每个固定 \(\omega\) 几乎处处相同。取可数稠密测试族并使用有限质量界，即可使零测例外集与 \(\omega\) 无关，得到流的恒等式。

	一般 Lipschitz \(f\) 可借助图嵌入 \(F(x)=(x,f(x))\) 化归这个光滑坐标情形。由\cref{g3:prop:rectifiable-lipschitz-pushforward,g3:cor:normal-pushforward-finite}，\(F_\#T\) 是正规可整流流，其取向由 \(\mathrm D_MF\) 给出；高度函数 \(q(x,s)=s\) 光滑。投影 \(\pi(x,s)=x\) 满足 \(\pi\circ F=\mathrm{id}\)，且 \(q\circ F=f\)。推前与边界相容，限制到相应水平下方也相容，故
	\[
		\pi_\#\langle F_\#T,q,t\rangle=\langle T,f,t\rangle.
	\]
	在 \(\nabla_Mf\ne0\) 的点，取切向正交基，使第一个向量沿 \(\nabla_Mf\)，其余向量位于水平切平面。图映射把它们送到 \((e_1,|\nabla_Mf|)\) 及 \((e_i,0)\)。高度切片后投影回去，正得到式\eqref{g3:eq:rectifiable-slice-representation}中的单位取向；图面积因子与高度余面积因子相消。这样便得到 Lipschitz 情形的同一公式。

	若 \(T\) 及 \(\partial T\) 都整数可整流，分别对它们使用这个表示，便知切片及其边界 \(-\langle\partial T,f,t\rangle\) 都整数可整流。截断流本身保持原重数，而其边界为两个整数可整流流之和。因此两者均为所述整数流。
\end{proof}

\begin{proposition}\label{g3:prop:iterated-current-slicing}
	设 \(T\in\mathcal N_k(\mathbb R^n)\)，\(1\le m\le k\)，\(f_i:\mathbb R^n\to\mathbb R\) 为 Lipschitz 函数。按 \(f_1,\ldots,f_m\) 的次序逐次切片，记为 \(\langle T,f,y\rangle\)，其中 \(f=(f_1,\ldots,f_m)\)、\(y\in\mathbb R^m\)。则它几乎处处为正规 \((k-m)\) 维流，且
	\[
		\partial\langle T,f,y\rangle=(-1)^m\langle\partial T,f,y\rangle,
		\qquad
		\int_{\mathbb R^m}\mathbf M(\langle T,f,y\rangle)\,\mathrm dy
		\le\prod_{i=1}^m\operatorname{Lip}(f_i)\,\mathbf M(T).
	\]
	若 \(T\) 整数，则这些切片也整数。若各 \(f_i\) 光滑，则对 \(\omega\in\mathcal D^{k-m}(\mathbb R^n)\) 和 \(g\in C_c^\infty(\mathbb R^m)\)，有
	\begin{equation}\label{g3:eq:iterated-slice-disintegration}
		\int_{\mathbb R^m}g(y)\langle T,f,y\rangle(\omega)\,\mathrm dy
		=T\bigl(g(f)\,\mathrm df_1\wedge\cdots\wedge\mathrm df_m\wedge\omega\bigr).
	\end{equation}
	特别地，当 \(m=k\)、\(p=(x_{i_1},\ldots,x_{i_k})\) 为有序坐标投影时，
	\[
		T(\varphi\,\mathrm dx_{i_1}\wedge\cdots\wedge\mathrm dx_{i_k})
		=\int_{\mathbb R^k}\langle T,p,y\rangle(\varphi)\,\mathrm dy
		\qquad(\varphi\in C_c^\infty(\mathbb R^n)).
	\]
\end{proposition}
\begin{proof}
	逐次应用\cref{g3:thm:normal-current-slicing}，并用 Fubini 定理选择共同的满测参数集。每施加一次边界就带来一个负号；每作一次质量积分就带来一个 Lipschitz 常数。整数性由\cref{g3:thm:integer-current-slicing}逐次保持。

	当 \(g(y)=g_1(y_1)\cdots g_m(y_m)\) 时，逐次使用式\eqref{g3:eq:current-slice-pairing}便得到式\eqref{g3:eq:iterated-slice-disintegration}，外微分因子按既定次序出现。有限个这种乘积的线性组合在紧参数盒上的连续函数中一致稠密；例如逐个坐标用分片线性插值，再将插值函数光滑逼近即可。两侧均有已经证明的质量积分界，所以公式推广到任意所述 \(g\)。最后取 \(g\) 在 \(p(\operatorname{spt}\varphi)\) 的邻域等于一，得到零维切片的配对公式。参数的可测性先对可数测试族成立，有限质量再将它推广到全部测试形式。
\end{proof}

一个切片仍保留取向。例如标准有向平面中的圆盘被 \(f(x,y)=x\) 切割，内部竖直弦的取向为 \(+e_y\)，因为 \((e_x\wedge e_y)\mathbin{\lrcorner}\mathrm dx=e_y\)。弦的上端点为正、下端点为负；这与圆周切片之间的负号恰好相配。

\begin{corollary}\label{g3:cor:coordinate-slice-atoms}
	设 \(T=\lBrack M,\tau,\theta\rBrack\in\mathcal N_k(\mathbb R^n)\) 为实重数可整流流，\(\theta>0\)。固定有序坐标投影 \(p=(x_{i_1},\ldots,x_{i_k})\)，记
	\[
		a_p(x)=\langle\tau(x),\mathrm dx_{i_1}\wedge\cdots\wedge\mathrm dx_{i_k}\rangle.
	\]
	则几乎所有零维切片有表示
	\begin{equation}\label{g3:eq:coordinate-slice-atoms}
		\langle T,p,y\rangle
		=\sum_{\substack{x\in M\cap p^{-1}(y)\\a_p(x)\ne0}}
		\theta(x)\operatorname{sgn}(a_p(x))\delta_x,
	\end{equation}
	各项绝对系数可求和，且
	\[
		\int_{\mathbb R^k}\mathbf M(\langle T,p,y\rangle)\,\mathrm dy
		=\int_M\theta|a_p|\,\mathrm d\mathcal H^k.
	\]
\end{corollary}
\begin{proof}
	逐次应用式\eqref{g3:eq:rectifiable-slice-representation}。在满秩点，连续的切向收缩最后留下标量 \(a_p\)，各步的余面积因子之积为切向 Jacobian \(J_Mp=|a_p|\)。秩亏部分在几乎所有最终纤维上为空，这由参数片上的\cref{g3:lem:coarea-critical-slices}给出。把带权余面积公式逐次积分，便得到末尾等式；其右端不超过 \(\mathbf M(T)\)，故几乎所有纤维上的绝对系数可求和。正重数的可求和性也说明非零原子至多可数。不同的纤维点对应不同原子，因此这里不发生同点抵消。
\end{proof}

\subsection{锥构造和填充估计}
\label{g3:subsec:2203:02}
给一个闭的有向曲面选定顶点，并把各点沿线段连向顶点，就得到一个候选填充。对于流，区间乘积和推前使这件事具有统一定义，也能直接算出质量。

\begin{proposition}\label{g3:prop:current-cone-filling}
	设 \(T\in\mathcal N_k(\mathbb R^n)\) 紧支，\(0\le k\le n\)，\(a\in\mathbb R^n\)。令
	\[
		h(t,x)=a+t(x-a),\qquad
		a*T=h_\#(\lBrack[0,1]\rBrack\times T).
	\]
	则 \(a*T\) 为正规 \((k+1)\) 维流，其支集包含于 \(a\) 与 \(\operatorname{spt}T\) 张成的闭线段并，并满足
	\begin{equation}\label{g3:eq:current-cone-boundary}
		\partial(a*T)=T-(c_a)_\#T-a*(\partial T),
	\end{equation}
	其中 \(c_a\) 是常值映射。质量估计为
	\begin{equation}\label{g3:eq:current-cone-mass}
		\mathbf M(a*T)\le\frac1{k+1}\int|x-a|\,\mathrm d\|T\|(x).
	\end{equation}
	若 \(T\) 整数，则 \(a*T\) 也整数。特别地，若 \(k\ge1\)、\(\partial T=0\) 且 \(\operatorname{spt}T\subset\overline B(a,R)\)，则 \(\partial(a*T)=T\)，且 \(\mathbf M(a*T)\le R\mathbf M(T)/(k+1)\)。若 \(k=0\)，同一填充结论须另要求 \(T(1)=0\)。
\end{proposition}
\begin{proof}
	区间取从零到一的取向。由\cref{g3:prop:current-homotopy}，其乘积的边界给出式\eqref{g3:eq:current-cone-boundary}。当 \(k\ge1\) 时，常值映射的 \(k\) 阶外幂微分为零，故 \((c_a)_\#T=0\)。当 \(k=0\) 时，\((c_a)_\#T=T(1)\delta_a\)；这里 \(T(1)\) 用在支集邻域等于一的紧支函数求值得到，与该函数的选择无关。

	在 \((t,x)\) 处，时间方向的微分为 \(x-a\)，空间方向的微分为乘以 \(t\)。于是对余质量至多一的 \((k+1)\) 形式，区间方向与空间 \(k\) 向量的配对绝对值不超过 \(|x-a|t^k\)。先对空间质量测度积分，再对 \(t\in[0,1]\) 积分，得到
	\[
		\mathbf M(a*T)\le\int_0^1\int|x-a|t^k\,\mathrm d\|T\|(x)\,\mathrm dt
		=\frac1{k+1}\int|x-a|\,\mathrm d\|T\|(x).
	\]
	边界公式同时给出边界的有限质量，所以锥为正规流。整数可整流流的区间乘积由原参数片乘区间表示，重数仍为原整数；其 Lipschitz 推前仍整数可整流。若 \(T\) 整数，边界公式中的各项也整数可整流，因而锥是整数流。最后代入 \(|x-a|\le R\)，并分别检查常值项，即得所述填充结论。
\end{proof}

零维的附加条件有具体含义。对 \(T=\delta_b-\delta_c\)，锥是从 \(a\) 到 \(b\) 的线段减去从 \(a\) 到 \(c\) 的线段，顶点贡献相消。若只有 \(T=\delta_b\)，其锥边界为 \(\delta_b-\delta_a\)。任何紧支一维流的边界总系数都为零，因为可用支集邻域等于一的测试函数计算；所以单个正点不能成为紧支流的边界。

当 \(T\) 是半径 \(R\) 的有向圆周且顶点在圆心时，锥是有向圆盘，质量估计成为 \(\pi R^2=(R/2)(2\pi R)\)。这个例子也说明，填充量的尺度比边界质量多一个长度因子。

\subsection{平坦范数与弱收敛}
\label{g3:subsec:2203:03}
两条距离很近但互不重合的曲线，其差的质量仍等于两条曲线的长度之和。然而它们之间的薄带面积很小。为了使这个几何差异能够趋零，我们同时允许留下少量流，并用低质量的高一维流填充其余部分。

\begin{definition}\label{g3:def:flat-norm}
	对 \(\mathbb R^n\) 中的 \(k\) 维流 \(T\)，定义实系数的\term[pingtan fanshu]{平坦范数}[flat norm]
	\begin{equation}\label{g3:eq:flat-norm}
		\mathcal F(T)=\inf\{\mathbf M(R)+\mathbf M(S):
		T=R+\partial S,\ \mathbf M(R)+\mathbf M(S)<\infty\},
	\end{equation}
	其中 \(R\) 为 \(k\) 维流，\(S\) 为 \((k+1)\) 维流；没有此类分解时取 \(+\infty\)。分解不要求整数系数。若 \(\mathcal F(T_j-T)\to0\)，称 \(T_j\) \term[pingtan shoulian]{平坦收敛}[flat convergence]于 \(T\)。
\end{definition}
\symidx{\(\mathcal F(T)\)：实系数流的平坦范数}

若 \(T\) 正规，则任意有限质量分解中的 \(R,S\) 自动正规：\(\partial S=T-R\) 有限质量，而 \(\partial R=\partial T\) 有限质量。这一事实使切片定理可以用于平坦范数中的每个分解。若特定论证要求 \(R,S\) 支于同一个闭集 \(K\)，则把对应下确界另记为 \(\mathcal F_K(T)\)；它可能更大，不能静默省略支集条件。

\begin{proposition}\label{g3:prop:flat-norm-duality}
	对 \(\mathbb R^n\) 中有限质量的 \(k\) 维流 \(T\)，有
	\begin{equation}\label{g3:eq:flat-dual-formula}
		\mathcal F(T)=\sup\{|T(\omega)|:
		\omega\in\mathcal D^k(\mathbb R^n),\quad
		\|\omega\|_\infty\le1,\ \|\mathrm d\omega\|_\infty\le1\},
	\end{equation}
	其中形式的范数均为余质量范数。\(\mathcal F\) 在有限质量流上是范数，并满足
	\[
		\mathcal F(T)\le\mathbf M(T),\qquad
		\mathcal F(\partial T)\le\mathcal F(T).
	\]
	第二式在边界具有有限平坦范数的流意义下成立。任意测试形式都有
	\begin{equation}\label{g3:eq:flat-controls-weak}
		|T(\omega)|\le\max\{\|\omega\|_\infty,\|\mathrm d\omega\|_\infty\}\,\mathcal F(T).
	\end{equation}
	因此平坦收敛蕴含流的弱收敛。
\end{proposition}
\begin{proof}
	若 \(T=R+\partial S\)，则
	\[
		|T(\omega)|\le\mathbf M(R)\|\omega\|_\infty
		+\mathbf M(S)\|\mathrm d\omega\|_\infty.
	\]
	对分解取下确界，得到式\eqref{g3:eq:flat-controls-weak}，并证明式\eqref{g3:eq:flat-dual-formula}中右端不超过左端。

	反向不等式使用\cref{g3:lem:bounded-functional-extension}。记右端上确界为 \(A\)，它不超过 \(\mathbf M(T)\)。在
	\[
		X=C_0(\mathbb R^n;\Lambda^k)\oplus C_0(\mathbb R^n;\Lambda^{k+1}),
		\qquad \|(\alpha,\beta)\|_X=\max\{\|\alpha\|_\infty,\|\beta\|_\infty\},
	\]
	的线性子空间 \(Y=\{(\omega,\mathrm d\omega):\omega\in\mathcal D^k\}\) 上，令 \(L(\omega,\mathrm d\omega)=T(\omega)\)。这个定义单值，且 \(\|L\|=A\)。保持范数延拓到 \(X\) 后，分别限制到两个直和项，得到连续形式上的泛函 \(R,S\)。由连续函数泛函的向量测度表示，它们是有限质量流。最大范数直和的对偶范数为两项范数之和：上界由三角不等式给出，下界可在两个分量中各选几乎取到范数的单位测试形式，并分别调整符号取得。因此
	\[
		\mathbf M(R)+\mathbf M(S)=A,\qquad
		T(\omega)=R(\omega)+S(\mathrm d\omega).
	\]
	这给出 \(T=R+\partial S\)，故 \(\mathcal F(T)\le A\)。这里使用的测度表示与\subsecref{g3:subsec:2101:02}相同，只把有限维向量的欧氏范数换成质量范数。

	齐次性由分解乘标量得到，三角不等式由两个近乎最优的分解相加得到。若 \(\mathcal F(T)=0\)，式\eqref{g3:eq:flat-controls-weak}说明它作用于每个测试形式都为零，故 \(T=0\)。取 \(R=T,S=0\)，得到 \(\mathcal F(T)\le\mathbf M(T)\)。最后由 \(T=R+\partial S\) 得 \(\partial T=\partial R\)，所以 \(\mathcal F(\partial T)\le\mathbf M(R)\)；取下确界便得边界估计。
\end{proof}

上述对偶表示还给出可测性：按紧支集穷尽，在两个范数都至多一的光滑形式中取可数的 \(C^1\) 稠密族，即可用同一可数族计算上确界。因此对切片而言，\(t\mapsto\mathcal F(\langle T,f,t\rangle)\) 是可测函数。

\begin{proposition}\label{g3:prop:flat-slicing-estimate}
	若 \(T\in\mathcal N_k(\mathbb R^n)\)，\(k\ge1\)，\(f:\mathbb R^n\to\mathbb R\) 为 \(L\)-Lipschitz 函数，则
	\begin{equation}\label{g3:eq:flat-slicing-estimate}
		\int_{\mathbb R}\mathcal F(\langle T,f,t\rangle)\,\mathrm dt
		\le L\mathcal F(T).
	\end{equation}
	若 \(T_j,T\) 都正规，且 \(\sum_j\mathcal F(T_j-T)<\infty\)，则对几乎所有 \(t\)，有 \(\mathcal F(\langle T_j,f,t\rangle-\langle T,f,t\rangle)\to0\)。所以平坦收敛于正规流的任意正规流序列，都可选出具有这一性质的子列。对固定的多参数切片，积分估计中的 \(L\) 换成各坐标 Lipschitz 常数之积。
\end{proposition}
\begin{proof}
	给定 \(\varepsilon>0\)，取 \(T=R+\partial S\)，使 \(\mathbf M(R)+\mathbf M(S)\le\mathcal F(T)+\varepsilon\)。如前所述，\(R,S\) 都正规。由切片的线性和边界公式，几乎处处有
	\[
		\langle T,f,t\rangle
		=\langle R,f,t\rangle-\partial\langle S,f,t\rangle.
	\]
	这本身就是一个平坦分解，故
	\[
		\mathcal F(\langle T,f,t\rangle)
		\le\mathbf M(\langle R,f,t\rangle)+\mathbf M(\langle S,f,t\rangle).
	\]
	积分后用\cref{g3:thm:normal-current-slicing}，得到上界 \(L(\mathcal F(T)+\varepsilon)\)；令 \(\varepsilon\downarrow0\) 即得估计。

	把估计用于 \(T_j-T\) 后求和，非负积分定理说明几乎所有 \(t\) 上相应平坦距离的级数收敛，因此各项趋零。若只有 \(\mathcal F(T_j-T)\to0\)，先选取距离小于 \(2^{-j}\) 的子列即可。多参数情形逐次积分并使用 Fubini 定理；只涉及可数个流和测试族，因此可把所有例外集并入一个零测集。
\end{proof}

这一估计说明，平坦控制可以随降维一起传递。闭包证明中的零维整数信息，就是通过反复切片从原序列传到几乎每个纤维上的。

\subsection{质量范数收敛、平坦收敛与弱收敛}
\label{g3:subsec:2203:04}
质量范数收敛指 \(\mathbf M(T_j-T)\to0\)，不是仅仅指两个数 \(\mathbf M(T_j)\to\mathbf M(T)\)。由上一小节有
\[
	\mathbf M(T_j-T)\longrightarrow0
	\quad\Longrightarrow\quad
	\mathcal F(T_j-T)\longrightarrow0
	\quad\Longrightarrow\quad T_j\rightharpoonup T.
\]
第一种收敛对不同支撑尤其敏感；后两种则允许一个有向几何对象通过很薄的区域移动。

\begin{example}\label{g3:exm:translated-segment-flat}
	在 \(\mathbb R^2\) 中，令 \(T_0\) 为从 \((0,0)\) 到 \((1,0)\) 的有向线段，\(T_\varepsilon\) 为从 \((0,\varepsilon)\) 到 \((1,\varepsilon)\) 的有向线段，\(\varepsilon>0\)。两条线段支撑互不相交，因而
	\[
		\mathbf M(T_\varepsilon-T_0)=2,
		\qquad \mathbf M(T_\varepsilon)=\mathbf M(T_0)=1.
	\]
	沿竖直平移作同伦，令 \(A_\varepsilon\) 为扫过的有向矩形流，\(B_\varepsilon\) 为两条端点轨迹组成的有向流。同伦公式给出
	\[
		T_\varepsilon-T_0=\partial A_\varepsilon+B_\varepsilon,
		\qquad \mathbf M(A_\varepsilon)=\varepsilon,
		\quad \mathbf M(B_\varepsilon)=2\varepsilon.
	\]
	所以 \(\mathcal F(T_\varepsilon-T_0)\le3\varepsilon\to0\)。总质量甚至保持不变，差的质量却不趋零。
\end{example}

\begin{example}\label{g3:exm:weak-not-flat-alternating-current}
	在环境 \(\mathbb R\) 中，将 \([0,1]\) 分成 \(2j\) 个等长小区间，交替赋予正、负取向，令
	\[
		T_j=\sum_{i=0}^{2j-1}(-1)^i
		\lBrack[i/(2j),(i+1)/(2j)]\rBrack.
	\]
	对光滑测试一形式 \(\varphi(x)\,\mathrm dx\)，将相邻两个区间配对，两个积分之差由 \(\varphi\) 在距离 \(1/(2j)\) 处的连续模控制。因此 \(T_j\rightharpoonup0\)。另一方面，\(\mathbf M(T_j)=1\)，而 \(\mathbb R\) 上没有非零二维流，故任何平坦分解中的 \(S\) 都为零，于是
	\[
		\mathcal F(T_j)=\mathbf M(T_j)=1.
	\]
	这些流都是整数流，支集也在同一个紧区间内；失去的条件是统一的边界质量界。内部 \(2j-1\) 个接点各有系数绝对值二，两个外端点各有系数绝对值一，故 \(\mathbf M(\partial T_j)=4j\to\infty\)。
\end{example}

弱收敛下，测试形式只能记录带符号的配对，质量因此可能因抵消而下降。平坦范数另外要求这些差异能够由低质量的本体与填充表示；质量范数则直接计算差的全部无符号质量。后面的紧性定理将在共同紧支和统一正规质量界下，把弱极限与平坦控制联系起来，同时证明极限保留整数可整流结构。
% 纲要条目排为编号小节标题，并进入目录与书签；说明段仍保留于纲要注释中。

% !TeX root = ../../Book3.tex
% 纲要标题：### 22.4 Federer–Fleming 变形定理
% 纲要位置：几何学计划纲要.md，第 3032 行。
% 正文状态：初稿；尚未正式审核。

% 纲要内容开始
% 
% * 立方网格与逐维径向投影
% * 质量控制和边界相容
% * 多面体逼近及误差填充
% * 网格尺度和统一常数
% 
% 纲要内容结束

\section{Federer–Fleming 变形定理}
\label{g3:sec:2204}
把一条曲线移到细网格上，误差可以由移动时扫过的带状区域填充。对高维流，困难在于同时控制面积和边界：径向投影在中心附近会放大质量，而投影到一个面上的流也未必具有常数重数。本节分别处理这两个问题。先选择质量分布较稀疏的投影中心，再用面内的 BV 估计把变动密度替换为常数。所证版本对应\cite[第29节，定理29.1、推论29.3及证明，第163--177页]{Simon1983}，以下展开所需构造与估计。

\subsection{立方网格与逐维径向投影}
\label{g3:subsec:2204:01}
固定环境 \(\mathbb R^n\) 和边长 \(h>0\)。以 \(h\mathbb Z^n\) 为顶点的闭立方体及其各维面组成网格；记 \(\mathcal Q_h^j\) 为所有 \(j\) 维面，\(\mathcal S_h^j\) 为它们的并。每个面选择一次坐标取向，\([F]\) 表示沿该取向积分的流；同伦中的 \([0,1]\) 同样表示取正向的区间流。这里网格的位置固定，随后改变的是立方体内部的投影中心。

取一个 \(j\) 维立方体 \(F\)，并在其中央半边长子立方体中选择 \(a_F\)。从 \(a_F\) 出发沿射线把 \(x\ne a_F\) 送到 \(\partial F\)，得到径向映射 \(r_F\)。它在边界上为恒等，在每个投向指定侧面的区域内有显式分式表达，因此
\[
	|r_F(x)-x|\le\sqrt j\,h,\qquad
	|\mathrm Dr_F(x)|\le \frac{C_jh}{|x-a_F|}.
\]
例如被击中的面为 \(x_i=b\) 时，
\[
	r_F(x)=a_F+\frac{b-(a_F)_i}{x_i-(a_F)_i}(x-a_F).
\]
在该区域内，分母的绝对值与 \(|x-a_F|\) 可用只依赖维数的常数比较；对上式求导即得估计。

\begin{lemma}\label{g3:lem:cubical-radial-retraction}
	记 \(\mathcal S_h^k\) 为标准立方网格 \(h\mathbb Z^n\) 中所有不超过 \(k\) 维闭面的并。
	设 \(0\le k<n\)、\(h>0\)，\(a\in(h/4,3h/4)^n\)。令
	\[
	\Sigma_a=\bigcup_{\substack{I\subset\{1,\ldots,n\}\\|I|=k+1}}
	\{x:x_i-a_i\in h\mathbb Z\text{ 对每个 }i\in I\}.
	\]
	则存在局部 Lipschitz 映射
	\(r:\mathbb R^n\setminus\Sigma_a\to\mathcal S_h^k\)，在 \(\mathcal S_h^k\) 上为恒等，每个立方体映入自身，并且
	\begin{equation}\label{g3:eq:cubical-radial-bound}
	|r(x)-x|\le\sqrt n\,h,\qquad
	\operatorname{Lip}_{\mathrm{loc}}r(x)\le
	\frac{C_{n,k}h}{\operatorname{dist}(x,\Sigma_a)}.
	\end{equation}
	局部 Lipschitz 常数的估计允许把右边乘一个固定的维数常数。
\end{lemma}
\begin{proof}
	在基本立方体 \([0,h]^n\) 内，先从 \(a\) 径向投到 \((n-1)\) 骨架，再在各 \((n-1)\) 面上，从 \(a\) 的正交投影继续向该面的边界投影。如此下降，直到到达 \(k\) 骨架。各步在较低维边界上固定不动，所以相邻面上的定义相容。平移同一构造便定义整个网格上的映射。

	每一步可能遇到的投影中心，以及此前径向映射下这些中心的原像，都包含于所列 \(\Sigma_a\)：在一个面内，中心的全部自由坐标等于对应 \(a_i\)；在下降到 \(k\) 维之前，至少还剩 \(k+1\) 个这样的坐标条件。因而在 \(\Sigma_a\) 外构造有定义。

	还须避免把逐步导数估计粗糙相乘而得到过强的奇性。下面对立方体维数归纳，证明总映射只有距离的一次倒数。对维数 \(k+1\)，这就是一次径向投影估计。假设更低维情况已证，令第一步把 \(x\) 送到面 \(F\) 内的 \(y\)，写
	\[
	y-a=\lambda(x-a),\qquad
	\lambda=\frac{|y-a|}{|x-a|},\qquad |y-a|\ge h/4.
	\]
	以 \(\pi_F\) 记到 \(F\) 的仿射正交投影，\(a_F=\pi_F(a)\)。面内余下步骤的奇性集是经过 \(a_F\) 的有限个坐标仿射子空间的并，记为 \(\Sigma_F\)。这些空间的法向圆柱均包含于基本立方体所用的奇性集 \(\Sigma\)。于是
	\[
	\begin{aligned}
	\operatorname{dist}(y,\Sigma_F)
	&=\lambda\operatorname{dist}(\pi_F(x),\Sigma_F)\\
	&\ge\lambda\operatorname{dist}(x,\Sigma).
	\end{aligned}
	\]
	归纳假设给出余下映射在 \(y\) 附近的 Lipschitz 常数至多
	\(Ch/\operatorname{dist}(y,\Sigma_F)\)。乘以第一步的 \(Ch/|x-a|\)，再用 \(|y-a|\ge h/4\)，便得到 \(Ch/\operatorname{dist}(x,\Sigma)\)。全局 \(\Sigma_a\) 包含该局部奇性集，所以也得到式\eqref{g3:eq:cubical-radial-bound}。在各分片交界处，映射连续且片数有限；沿线段逐片求和给出同样的局部 Lipschitz 估计。各步留在原立方体中，故位移不超过其直径。
\end{proof}

奇性集不是任意选取的。下一条估计使流在其小邻域内的质量衰减得足够快，抵消径向投影的导数增长。

\begin{lemma}\label{g3:lem:deformation-good-center}
	设 \(0\le k<n\)，\(\mu_1,\mu_2\) 为 \(\mathbb R^n\) 上有限非负 Radon 测度，\(h>0\)。存在 \(a\in(h/4,3h/4)^n\)，使 \(\Sigma_a=\bigcup_{|I|=k+1}\{x:x_i-a_i\in h\mathbb Z\ (i\in I)\}\) 满足，对 \(i=1,2\) 及 \(0<\rho<h\)，
	\begin{equation}\label{g3:eq:good-center-tube}
	\mu_i\{x:\operatorname{dist}(x,\Sigma_a)<\rho\}
	\le C_{n,k}(\rho/h)^{k+1}\mu_i(\mathbb R^n).
	\end{equation}
\end{lemma}
\begin{proof}
	零测度无须限制中心。对非零 \(\mu\) 及一个大小为 \(k+1\) 的坐标集合 \(I\)，将 \(\mu\) 按这些坐标投影，再模去 \(h\mathbb Z^{k+1}\)，得到边长 \(h\) 的平坦环面上的测度 \(\nu_I\)。其总质量为 \(\mu(\mathbb R^n)\)。将满足
	\[
	\nu_I(B_\rho(b))>A(\rho/h)^{k+1}\mu(\mathbb R^n)
	\quad\text{对某个 }0<\rho<h/10
	\]
	的 \(b\) 称为坏点。五倍球覆盖引理给出两两不交的坏球，其五倍球覆盖所有坏点。把各球的不等式相加，得到坏点集合的 \((k+1)\) 维体积至多 \(C_{n,k}h^{k+1}/A\)。这里可先处理有限子族，再用外正则性与可数覆盖取极限；环面上的小球在有限张坐标图内就是 Euclid 球。

	将坏点集合沿其余坐标拉回中央子立方体，所占 \(n\) 维体积至多 \(C_{n,k}h^n/A\)。对有限个 \(I\) 和两项测度取并，选择仅依赖 \(n,k\) 的充分大 \(A\)，坏点总体积便小于中央子立方体的体积。取其外的 \(a\)。距 \(\Sigma_a\) 小于 \(\rho\) 的点必属于至少一个坐标投影逆像球，故把相应良好估计相加即得 \(\rho<h/10\) 时的结论。其余半径用总质量估计，并增大维数常数即可。
\end{proof}

\subsection{质量控制和边界相容}
\label{g3:subsec:2204:02}
将上条引理同时用于 \(\|T\|\) 与 \(\|\partial T\|\)。选择中心时同时照顾这两项测度，正是后续边界估计仍能成立的原因。

\begin{proposition}\label{g3:prop:deformation-to-skeleton}
	下式 \(\mathcal S_h^k\) 表示网格 \(h\mathbb Z^n\) 的 \(k\) 维骨架。
	设 \(0\le k<n\)，\(T\in\mathcal N_k(\mathbb R^n)\) 紧支，\(h>0\)。存在紧支正规流 \(Q,R_0,S_0\)，其维数依次为 \(k,k+1,k\)，满足
	\[
	T-Q=\partial R_0+S_0,\qquad
	\operatorname{spt}Q\subset\mathcal S_h^k,
	\]
	以及
	\begin{equation}\label{g3:eq:skeleton-deformation-bounds}
	\begin{aligned}
	\mathbf M(Q)&\le C\mathbf M(T),&
	\mathbf M(\partial Q)&\le C\mathbf M(\partial T),\\
	\mathbf M(R_0)&\le Ch\mathbf M(T),&
	\mathbf M(S_0)&\le Ch\mathbf M(\partial T).
	\end{aligned}
	\end{equation}
	\(Q,R_0\) 支于 \(\operatorname{spt}T\) 的 \(\sqrt n h\) 邻域，\(\partial Q,S_0\) 支于 \(\operatorname{spt}\partial T\) 的同样邻域。若 \(T\) 为整数流，则 \(Q,R_0,S_0\) 均为整数流。常数 \(C\) 只依赖 \(n,k\)。
\end{proposition}
\begin{proof}
	记 \(d(x)=\operatorname{dist}(x,\Sigma_a)\)，其中 \(a\) 按\cref{g3:lem:deformation-good-center}选取；令 \(r\) 为相应径向映射。对 \(\mu=\|T\|\) 或 \(\|\partial T\|\)，式\eqref{g3:eq:good-center-tube}与二进环带求和给出，对 \(0\le q\le k\)，
	\begin{equation}\label{g3:eq:radial-weight-tail}
	\int_{\{0<d<\rho\}}(h/d)^q\,\mathrm d\mu
	\le C(\rho/h)^{k+1-q}\mu(\mathbb R^n).
	\end{equation}
	事实上，第 \(j\) 个环带 \(2^{-j-1}\rho\le d<2^{-j}\rho\) 的贡献至多
	\(C(h/(2^{-j-1}\rho))^q(2^{-j}\rho/h)^{k+1}\mu(\mathbb R^n)\)，其比值为 \(2^{-(k+1-q)}<1\)。特别地，\(\mu(\Sigma_a)=0\)，而完整加权积分至多为 \(C\mu(\mathbb R^n)\)。

	先假设 \(k\ge1\)。切片质量估计\eqref{g3:eq:normal-slice-mass}允许在每个二进区间 \((\rho/2,\rho)\) 内选取良好半径 \(t\)，使
	\[
	\mathbf M\langle T,d,t\rangle
	\le C\rho^k h^{-k-1}\mathbf M(T).
	\]
	同时避开所有限制、切片与整数性结论的零测例外集，取这样的一列 \(t_j\downarrow0\)。写
	\(T_j=T\llcorner\{d>t_j\}\)，由切片定义与 \(\partial T\) 的限制公式，
	\[
	\partial T_j=(\partial T)\llcorner\{d>t_j\}-\langle T,d,t_j\rangle.
	\]
	将 \(r\) 和线性同伦
	\(H(s,x)=(1-s)x+sr(x)\) 用于这些支集远离 \(\Sigma_a\) 的正规流。推前只需要映射在支集邻域 Lipschitz，可先延拓再使用\cref{g3:prop:current-homotopy}。

	位移至多为 \(\sqrt n h\)，故推前的 \(q\) 维质量放大因子至多为 \(C(h/d)^q\)，而同伦作用于 \(q\) 维流时的质量放大因子至多为 \(Ch(h/d)^q\)。在 \(d\) 远离零的部分可以把 \(h/d\) 换成 \(1+h/d\)，不改变以下估计。式\eqref{g3:eq:radial-weight-tail}因此说明，下列三列均在质量范数中收敛：
	\[
	\begin{aligned}
	Q&=\lim_j r_\#T_j,\\
	R_0&=-\lim_j H_\#([0,1]\times T_j),\\
	S_0&=-\lim_j H_\#([0,1]\times
	((\partial T)\llcorner\{d>t_j\})).
	\end{aligned}
	\]
	例如两个不同删管尺度之间的 \(r_\#T\) 差，由 \(q=k\) 的尾估计控制为 \(C\rho\mathbf M(T)/h\)；相应同伦差至多为 \(C\rho\mathbf M(T)\)。有限质量流空间是有限维向量测度的完备空间，所以这些极限存在。

	还须处理删管所产生的边界。对上述切片，由 \(d=t_j\) 上的局部 Lipschitz 界，
	\[
	\begin{aligned}
	\mathbf M(r_\#\langle T,d,t_j\rangle)
	&\le C(t_j/h^2)\mathbf M(T)\longrightarrow0,\\
	\mathbf M(H_\#([0,1]\times\langle T,d,t_j\rangle))
	&\le C(t_j/h)\mathbf M(T)\longrightarrow0.
	\end{aligned}
	\]
	第一式使用切片维数 \(k-1\)，第二式再乘位移尺度 \(h\)。同伦恒等式取极限便得 \(T-Q=\partial R_0+S_0\)；边界推前公式同时给出
	\[
	\partial Q=\lim_j r_\#((\partial T)\llcorner\{d>t_j\}).
	\]
	用完整加权积分分别估计各极限，得到式\eqref{g3:eq:skeleton-deformation-bounds}。\(\partial S_0=\partial T-\partial Q\)，且 \(\partial R_0=T-Q-S_0\)，所以三个流均正规。支集控制来自每个点和它的同伦轨迹都留在原立方体内。

	若 \(T\) 为整数流，良好限制和切片由\cref{g3:thm:integer-current-slicing}仍为整数流；Lipschitz 推前及乘积保持整数可整流性。上述三个质量范数极限由\cref{g3:lem:integer-rectifiable-mass-closure}仍整数可整流，边界公式再给出它们的边界也整数可整流。这里用到的是质量范数下的闭合性。

	当 \(k=0\) 时，\(\partial T=0\)，不需要切片项。对有限测度的限制直接定义 \(Q=r_\#T\) 与 \(R_0=-H_\#([0,1]\times T)\)，式\eqref{g3:eq:radial-weight-tail}在 \(q=0\) 时保证删管极限存在，零维同伦公式给出 \(T-Q=\partial R_0\)。取 \(S_0=0\) 即得全部结论；整数零维流是有限个带整数系数点，其每个点只沿一条线段移动。
\end{proof}

\subsection{多面体逼近及误差填充}
\label{g3:subsec:2204:03}
支于 \(k\) 骨架还不等于多面体流。一个面内仍可能有随位置变化的密度；其变化由面内的边界质量控制。下面先建立所需的 BV 估计。

\begin{lemma}\label{g3:lem:cube-median-bv}
	设 \(k\ge1\)，\(F\subset\mathbb R^k\) 为边长 \(h\) 的开立方体，\(u\in BV(F)\)。存在实数中位数 \(b\)，满足
	\[
	\mathcal L^k\{u\le b\}\ge|F|/2,\qquad
	\mathcal L^k\{u\ge b\}\ge|F|/2,
	\]
	并有
	\begin{equation}\label{g3:eq:cube-median-estimates}
	\begin{aligned}
	|b|h^k&\le2\int_F|u|,\\
	\int_F|u-b|&\le C_kh|Du|(F),\\
	|D((u-b)\chi_F)|(\mathbb R^k)&\le C_k|Du|(F).
	\end{aligned}
	\end{equation}
	最后一行将 \(u-b\) 在 \(F\) 外补零。若 \(u\) 几乎处处取整数值，\(b\) 可选为整数。
\end{lemma}
\begin{proof}
	分布函数给出中位数；整数值情形选取累计分布首次达到一半的整数即可。若 \(b>0\)，至少一半点上 \(|u|\ge b\)；若 \(b<0\)，至少一半点上 \(|u|\ge-b\)。这证明第一行。

	对光滑 \(v\)，在立方体内逐个坐标连接两点，对每条坐标线用微积分基本定理，再对两点积分，得到
	\[
	\int_F|v-v_F|\le C_kh\int_F|\nabla v|,
	\qquad v_F=|F|^{-1}\int_Fv.
	\]
	例如每次只替换一个坐标，其差由该坐标方向导数在长度 \(h\) 的区间上积分控制；对其余坐标积分并把 \(k\) 项相加，即得此式。该式也适用于 BV 函数：以立方体中心为中心先作比率 \(1-\delta\) 的内缩，再作小尺度卷积，得到定义于闭立方体邻域的光滑函数。内缩后的变差不超过
	\((1-\delta)^{1-k}|Du|(F)\)，函数在 \(L^1\) 中趋于 \(u\)；随后选取卷积尺度并用下半连续性，就得到保持总变差的光滑逼近。

	中位数使函数 \(c\mapsto\int_F|u-c|\) 取得最小值。具体地，若 \(c>b\)，在 \(u\le b\) 的至少一半点上绝对值增加 \(c-b\)，其余点上至多减少 \(c-b\)；\(c<b\) 同理。因此
	\(\int_F|u-b|\le\int_F|u-u_F|\)，证明第二行。

	对定义于闭立方体邻域的光滑 \(v\)，沿垂直每个侧面的线段积分，得到
	\[
	\int_{\partial F}|v|\,\mathrm d\mathcal H^{k-1}
	\le C_k\left(h^{-1}\int_F|v|+\int_F|\nabla v|\right).
	\]
	一条线段上的估计为
	\(|v(0)|\le h^{-1}\int_0^h|v(t)|\,\mathrm dt+\int_0^h|v'(t)|\,\mathrm dt\)，另一端相同。补零的分布导数包含内域梯度和边界跳跃，所以
	\[
	|D(v\chi_F)|(\mathbb R^k)
	\le C_k\left(h^{-1}\int_F|v|+\int_F|\nabla v|\right).
	\]
	用上一段构造的严格光滑逼近处理 \(u-b\)，在全空间使用变差下半连续性，得到同一补零估计。结合第二行即为第三行。
\end{proof}

\termidx[Federer Fleming bianxing dingli]{Federer--Fleming 变形定理}
\foreignidx{Federer--Fleming deformation theorem}
\begin{theorem}[Federer--Fleming 变形定理]\label{g3:thm:federer-fleming-deformation}
	设 \(0\le k\le n\)，\(T\in\mathcal N_k(\mathbb R^n)\) 紧支，\(h>0\)。记 \(\mathcal Q_h^k\) 为标准立方网格的定向 \(k\) 维面族，\([F]\) 为对应积分流。在固定网格 \(h\mathbb Z^n\) 上，存在有限多面体流
	\[
	P=\sum_{F\in\mathcal Q_h^k}b_F[F]
	\]
	及紧支正规流 \(R\in\mathcal N_{k+1}(\mathbb R^n)\)、\(S\in\mathcal N_k(\mathbb R^n)\)，使
	\begin{equation}\label{g3:eq:deformation-decomposition}
	T-P=\partial R+S
	\end{equation}
	且
	\begin{equation}\label{g3:eq:deformation-mass-bounds}
	\begin{aligned}
	\mathbf M(P)&\le C\mathbf M(T),&
	\mathbf M(\partial P)&\le C\mathbf M(\partial T),\\
	\mathbf M(R)&\le Ch\mathbf M(T),&
	\mathbf M(S)&\le Ch\mathbf M(\partial T).
	\end{aligned}
	\end{equation}
	常数 \(C=C(n,k)\) 与 \(T,h\) 无关；\(P,R\) 支于 \(\operatorname{spt}T\) 的 \(2\sqrt n h\) 邻域，\(\partial P,S\) 支于 \(\operatorname{spt}\partial T\) 的同样邻域。若 \(T\in\mathcal I_k(\mathbb R^n)\)，则所有 \(b_F\in\mathbb Z\)，且 \(R,S\) 也是整数流。特别地，\(\partial T=0\) 时可取 \(S=0\) 并有 \(\partial P=0\)。
\end{theorem}
\begin{proof}
	当 \(k<n\) 时，先用\cref{g3:prop:deformation-to-skeleton}得到 \(Q,R_0,S_0\)。当 \(k=n\) 时，取 \(Q=T\)、\(R_0=S_0=0\)；全空间本来就是 \(n\) 骨架。若 \(k=0\)，\(Q\) 已经是顶点上的有限线性组合，直接令 \(P=Q\)、\(R=R_0\)、\(S=0\) 即可。以下设 \(k\ge1\)。

	对每个 \(k\) 面 \(F\)，在其相对内部 \(F^\circ\) 的小邻域中，\(Q\) 只支于该面的仿射平面。由\cref{g3:prop:normal-plane-bv}，
	\[
	Q\llcorner F^\circ=u_F[F],\qquad u_F\in BV(F^\circ),
	\]
	并且
	\[
	\int_{F^\circ}|u_F|=\|Q\|(F^\circ),\qquad
	|Du_F|(F^\circ)\le\|\partial Q\|(F^\circ).
	\]
	\cref{g3:lem:normal-current-no-lower-skeleton}说明 \(Q\) 不给低维骨架质量，故这些面内表示相加就是全部 \(Q\)。只有有限个面与 \(\operatorname{spt}Q\) 相交。

	在每个面选取\cref{g3:lem:cube-median-bv}的中位数 \(b_F\)，令 \(P=\sum_Fb_F[F]\)。中位数第一条估计给出
	\[
	\mathbf M(P)=\sum_F|b_F|h^k\le2\mathbf M(Q).
	\]
	将另外两条估计求和，得到
	\begin{equation}\label{g3:eq:skeleton-median-error}
	\mathbf M(Q-P)\le Ch\mathbf M(\partial Q),\qquad
	\mathbf M(\partial(Q-P))\le C\mathbf M(\partial Q).
	\end{equation}
	在第二式中，每个面上的 \((u_F-b_F)[F]\) 都先补零，故它在面边界上的跳跃也已包括；相邻面抵消只能减小求和后的质量。由 \(\partial P=\partial Q-\partial(Q-P)\)，再得 \(\mathbf M(\partial P)\le C\mathbf M(\partial Q)\)。

	现在取
	\[
	R=R_0,\qquad S=S_0+(Q-P).
	\]
	代入 \(T-Q=\partial R_0+S_0\) 即得分解，前述估计与式\eqref{g3:eq:skeleton-deformation-bounds}给出全部质量界。若面内 \(|Du_F|\) 为零，\(u_F\) 几乎处处为常数，\(Q-P\) 在该面上为零。因此 \(Q-P\) 所占的面必接触 \(\operatorname{spt}\partial Q\)，其每一点距此支集至多为一个面直径。这与先前支集控制相加，给出 \(S,\partial P\) 的 \(2\sqrt n h\) 邻域界；\(P,R\) 同理。

	若 \(T\) 为整数流，\(Q,R_0,S_0\) 为整数流，\(u_F\) 几乎处处为整数，可以选取整数中位数。所以 \(P\) 及其多面体边界都是整数可整流流，\(Q-P\) 及其边界也整数可整流，所构造的 \(R,S\) 是整数流。\(\partial T=0\) 时，式\eqref{g3:eq:deformation-mass-bounds}迫使 \(S=0\)、\(\partial P=0\)。
\end{proof}

该分解区分两种误差：\(R\) 记录流本身移动时扫过的 \((k+1)\) 维区域，\(S\) 记录边界移动和面内密度调整造成的 \(k\) 维误差。闭流的边界没有移动，第二种误差完全消失。对一般有边界流，\(S\) 的质量随网格尺度成正比地趋零。

\subsection{网格尺度和统一常数}
\label{g3:subsec:2204:04}
以上证明中，所有选点坏集比例、坐标投影片数和中位数估计的常数都只依赖维数。也可先在单位网格完成构造，再用伸缩 \(d_h(x)=hx\) 转回一般网格：\(j\) 维流的质量乘以 \(h^j\)，所以 \(R\) 比 \(T\) 多一个 \(h\)，\(S\) 比 \(\partial T\) 多一个 \(h\)。这说明误差估计中的尺度幂次来自维数，而不是逐个流另选的常数。

\begin{corollary}\label{g3:cor:uniform-polyhedral-flat-approximation}
	设 \(K\subset\mathbb R^n\) 紧，\(A<\infty\)。对所有支于 \(K\) 且
	\(\mathbf M(T)+\mathbf M(\partial T)\le A\) 的正规 \(k\) 流，在每个固定网格尺度 \(h>0\) 上均可选取多面体流 \(P_h(T)\)，使
	\[
	\mathcal F(T-P_h(T))\le C_{n,k}hA.
	\]
	当 \(h\) 固定时，这些多面体流属于同一个有限维向量空间，其系数一致有界。
\end{corollary}
\begin{proof}
	将\eqref{g3:eq:deformation-decomposition}代入平坦范数的定义，得到
	\(\mathcal F(T-P)\le\mathbf M(R)+\mathbf M(S)\le ChA\)。多面体支集包含于固定紧邻域 \(K_{2\sqrt n h}=\{x:\operatorname{dist}(x,K)\le2\sqrt n h\}\)，只涉及其中有限个网格面。各面内部不交，故
	\(\mathbf M(P)=h^k\sum_F|b_F|\le CA\)，从而每个系数的绝对值至多为 \(CAh^{-k}\)。
\end{proof}

这条推论把一个无限维流族，在任意给定误差下放入一个有限维有界集合的邻域。下一节的紧性证明将先固定 \(h\) 处理这些有限个系数，再令 \(h\) 趋零。

整数系数还给出一个不同用途：非零网格流至少占据一整份面质量。因此质量小于 \(h^k\) 的整数网格 \(k\) 流只能为零。

\begin{corollary}\label{g3:cor:integral-cycle-isoperimetric}
	设 \(1\le k<n\)，\(T\in\mathcal I_k(\mathbb R^n)\) 紧支且 \(\partial T=0\)。存在紧支整数 \((k+1)\) 流 \(V\)，满足
	\[
	\partial V=T,\qquad
	\mathbf M(V)\le C_{n,k}\mathbf M(T)^{(k+1)/k}.
	\]
	若 \(T\ne0\)，可使 \(V\) 支于 \(\operatorname{spt}T\) 的
	\(C_{n,k}\mathbf M(T)^{1/k}\) 邻域。
\end{corollary}
\begin{proof}
	\(T=0\) 时取 \(V=0\)。否则记变形定理中控制 \(\mathbf M(P)\) 的常数为 \(C_0\)，选择
	\(h=(2C_0\mathbf M(T))^{1/k}\)。由于 \(P\) 的系数为整数，若 \(P\ne0\)，则
	\[
	h^k\le\mathbf M(P)\le C_0\mathbf M(T)=h^k/2,
	\]
	矛盾。故 \(P=0\)。又因 \(\partial T=0\)，\(S=0\)，分解变为 \(T=\partial R\)。取 \(V=R\)，质量和支集界由变形定理立即得到。
\end{proof}

这里得到的是维数常数控制的填充估计；第21章为余维一集合证明的尖锐 Euclid 等周不等式还确定了最佳常数。两种估计分别服务于几何极值和一般整数流的填充。
% 纲要条目排为编号小节标题，并进入目录与书签；说明段仍保留于纲要注释中。

% !TeX root = ../../Book3.tex
% 纲要标题：### 22.5 整数流的闭包与紧性
% 纲要位置：几何学计划纲要.md，第 3039 行。
% 正文状态：初稿；尚未正式审核。

% 纲要内容开始
% 
% * 正规流的弱紧性
% 
%   在共同紧支和质量、边界质量界下，先获得对测试形式收敛的子列。
% 
% * 变形估计与统一平坦控制
% 
%   利用多面体逼近把弱序列转化为可控制的几何误差。
% 
% * 切片及维数归纳
% 
%   证明几乎所有切片保留相容的整数信息，逐维恢复整性。
% 
% * 极限的可整流性
% 
%   此处单列闭包定理真正需要的切向结构论证，不从弱收敛直接断言。
% 
% * 整数重数的保持
% 
%   区分可整流性与重数属于整数这两项结论。
% 
% * 边界相容及下半连续
% 
%   证明极限边界正确，并核验质量下半连续。
% 
% * 整数流的整体与局部紧性
% 
%   先完成固定紧支版本，再说明局部质量界与质量逃逸的区别。
% 
% ---
% 
% 纲要内容结束

\section{整数流的闭包与紧性}
\label{g3:sec:2205}

一列定向曲面可以发生折叠与抵消。质量界允许我们对测试形式取极限，变形定理又把这种收敛加强为平坦收敛；要使极限仍能作为整数 Plateau 问题的竞争者，还须恢复它的切平面和整数重数。本节从零维切片入手：切片中的点带整数权，正规性控制这些点随切片参数的移动，进而把它们组织成可数个 Lipschitz 图像。

\ProseParagraphBreak
以下先在 \(\mathbb R^n\) 中处理共同紧支的流，并记
\[
	\mathbf N(T)=\mathbf M(T)+\mathbf M(\partial T).
\]
切片与 BV 相结合的证明方法可参看 \cite[第7节，引理7.3、定理7.4；第8节，定理8.1与8.5，第42--49页]{AmbrosioKirchheim2000Currents}。这里把所需论证写成 Euclid 空间中的坐标测度与 Lipschitz 图像构造；所用测度、BV 和面积公式均已在前两章建立。

\subsection{正规流的弱紧性}
\label{g3:subsec:2205:01}
先保留能够直接从对偶性得到的结论。它适用于实系数正规流，因此尚未涉及整数重数。

\begin{theorem}\label{g3:thm:normal-current-weak-compactness}
	设 \(0\le k\le n\)，\(K\subset\mathbb R^n\) 紧，\(T_j\) 为正规 \(k\) 流，且
	\[
		\operatorname{spt}T_j\subset K,\qquad
		\sup_j\mathbf N(T_j)\le A<\infty.
	\]
	则存在子列及支于 \(K\) 的正规 \(k\) 流 \(T\)，使 \(T_j\rightharpoonup T\)、\(\partial T_j\rightharpoonup\partial T\)。对每个开集 \(U\subset\mathbb R^n\)，有
	\begin{equation}\label{g3:eq:normal-current-limit-mass}
		\|T\|(U)\le\liminf_j\|T_j\|(U),\qquad
		\|\partial T\|(U)\le\liminf_j\|\partial T_j\|(U).
	\end{equation}
\end{theorem}
\begin{proof}
	由\cref{g3:prop:current-mass-representation}，每个 \(T_j\) 是有限维空间 \(\Lambda_k\mathbb R^n\) 值的测度，其全变差按质量范数计量。取一个包含 \(K\) 的紧球，球上连续 \(k\) 形式的一致范数空间可分：在有限个坐标系数中用网格上的分片仿射函数逼近，再把网格与函数值取为有理数，就得到可数稠密族。各 \(T_j\) 在每个族中元素上的值一致有界。逐个取收敛子列再取对角线，所得子列在这一可数族上全部收敛。

	对任意连续形式 \(\omega\)，以该族中的 \(\omega_\ell\) 一致逼近。估计
	\[
		|T_j(\omega)-T_j(\omega_\ell)|
		\le A\|\omega-\omega_\ell\|_\infty
	\]
	与 \(j\) 无关，因此极限延拓为连续线性泛函。有限维分量上的 Radon 测度表示给出质量有限的流 \(T\)。若 \(\omega\) 支于 \(K\) 外，则全部 \(T_j(\omega)=0\)，所以 \(T(\omega)=0\)，即 \(\operatorname{spt}T\subset K\)。

	对光滑紧支 \((k-1)\) 形式 \(\eta\)，有
	\[
		\partial T(\eta)=T(\mathrm d\eta)
		=\lim_jT_j(\mathrm d\eta)=\lim_j\partial T_j(\eta),
		\qquad |\partial T(\eta)|\le A\|\eta\|_\infty.
	\]
	故 \(\partial T\) 也有有限质量，且边界沿同一子列收敛。最后对支于任意开集 \(U\) 的单位余质量测试形式取上确界，得到式\eqref{g3:eq:normal-current-limit-mass}；这也正是\cref{g3:prop:current-mass-lsc}的应用。\(k=0\) 时边界按约定为零，证明中的边界部分省去。
\end{proof}

这里的收敛允许质量损失。即使每个 \(T_j\) 的质量都等于一，极限质量也可以更小；几何上，相反取向的部分可以在靠近时相消。

\subsection{变形估计与统一平坦控制}
\label{g3:subsec:2205:02}
变形定理的作用，是把无限维的流族压缩到固定网格上的有限维系数空间。网格的径向中心可以随流改变，最终使用的立方骨架则保持不变。

\begin{proposition}\label{g3:prop:normal-bounded-flat-precompactness}
	设 \(K\subset\mathbb R^n\) 紧，\(0\le k\le n\)，并令
	\[
		\mathcal A_{K,A}
		=\{T:\ T\text{ 是正规 }k\text{ 流},\
		\operatorname{spt}T\subset K,\ \mathbf N(T)\le A\}.
	\]
	这个集合在平坦范数下序列紧。若 \(T_j,T\in\mathcal A_{K,A}\)，则
	\[
		T_j\rightharpoonup T
		\quad\Longleftrightarrow\quad
		\mathcal F(T_j-T)\longrightarrow0.
	\]
\end{proposition}
\begin{proof}
	固定 \(0<h<1\)。\cref{g3:thm:federer-fleming-deformation}给出
	\[
		T_j-P_j=\partial R_j+S_j,
		\quad
		\mathbf M(P_j)+\mathbf M(\partial P_j)\le C A,
		\quad
		\mathbf M(R_j)+\mathbf M(S_j)\le ChA,
	\]
	其中 \(P_j\) 是固定 \(h\) 网格上的多面体流，\(P_j,R_j,S_j\) 支于 \(K\) 的同一个 \(C_nh\) 邻域。只有有限个网格面与这个邻域相交。写
	\[
		P_j=\sum_{F}a_{j,F}\lBrack F\rBrack.
	\]
	互不相交的面内域各有体积 \(h^k\)，故
	\(\sum_F|a_{j,F}|h^k=\mathbf M(P_j)\le CA\)。系数落在有限维有界集内，因而可取子列，使 \(P_j\) 在质量范数中收敛于某个多面体流 \(P\)。

	现在假设 \(T_j\rightharpoonup T\)。除了上述质量界，还有
	\[
		\partial R_j=T_j-P_j-S_j,\qquad
		\partial S_j=\partial T_j-\partial P_j.
	\]
	因此 \(R_j,S_j\) 的边界质量也一致有界。对它们应用\cref{g3:thm:normal-current-weak-compactness}，再取子列得到 \(R_j\rightharpoonup R\)、\(S_j\rightharpoonup S\)，且
	\[
		T-P=\partial R+S,\qquad
		\mathbf M(R)+\mathbf M(S)\le ChA.
	\]
	利用这两个分解，得到
	\[
		\mathcal F(T_j-T)
		\le\mathbf M(P_j-P)
		+\mathbf M(R_j)+\mathbf M(S_j)+\mathbf M(R)+\mathbf M(S),
	\]
	所以沿所取子列有 \(\limsup_j\mathcal F(T_j-T)\le2ChA\)。

	若原序列不平坦收敛于 \(T\)，可先取一个满足 \(\mathcal F(T_j-T)\ge\varepsilon\) 的子列，再选 \(h\) 使 \(2ChA<\varepsilon\)。刚才的论证又从它抽出一个违反此下界的子列，矛盾。因此整个序列平坦收敛。反向由
	\[
		|(T_j-T)(\omega)|
		\le\max\{\|\omega\|_\infty,\|\mathrm d\omega\|_\infty\}
		\mathcal F(T_j-T)
	\]
	立即得到。任取 \(\mathcal A_{K,A}\) 中的序列，先用弱紧性取子列，再使用已经证明的等价性，就得到平坦序列紧性。
\end{proof}

当输入是整数流时，多面体系数也是整数，质量界甚至使固定网格上的可能系数组合只有有限个。不过上面的证明容许实系数，这一点使我们能够先控制极限，再判断它的整数性。

\subsection{切片及维数归纳}
\label{g3:subsec:2205:03}
零维整数流只有有限个点和有限个整数权，这是维数归纳的起点。

\begin{lemma}\label{g3:lem:integer-zero-current-closure}
	设 \(K\subset\mathbb R^n\) 紧，\(A_j\) 是支于 \(K\) 的整数零流，\(\sup_j\mathbf M(A_j)<\infty\)。若 \(A_j\rightharpoonup A\)，则 \(A\) 仍是整数零流。
\end{lemma}
\begin{proof}
	把每个非零整数系数拆成若干个 \(+1\) 或 \(-1\)，可写成
	\[
		A_j=\sum_{\ell=1}^{N_j}\varepsilon_{j,\ell}\delta_{x_{j,\ell}},
		\qquad \varepsilon_{j,\ell}\in\{-1,1\},\quad
		N_j=\mathbf M(A_j)\le C.
	\]
	先取子列固定 \(N_j=N\) 和有限个符号，再由 \(K^N\) 的紧性取子列，使 \(x_{j,\ell}\to x_\ell\)。于是
	\(A=\sum_{\ell=1}^N\varepsilon_\ell\delta_{x_\ell}\)。
	重合点的整数系数相加，系数为零的点删去，就得到整数零流的表示。这里使用子列只为识别已经给定的弱极限。
\end{proof}

\begin{lemma}\label{g3:lem:slice-closure-subsequence}
	设 \(1\le k\le n\)，\(T_j,T\) 为支于同一紧集 \(K\) 的正规 \(k\) 流，
	\[
		T_j\rightharpoonup T,\qquad \sup_j\mathbf N(T_j)<\infty.
	\]
	对任意固定 Lipschitz 函数 \(f:\mathbb R^n\to\mathbb R\)，几乎每个 \(t\) 都有一个可依赖于 \(t\) 的子列，使
	\[
		\langle T_j,f,t\rangle\rightharpoonup\langle T,f,t\rangle,
		\qquad
		\sup_j\mathbf N(\langle T_j,f,t\rangle)<\infty.
	\]
	若这些 \(T_j\) 为整数流，且整数流闭包结论已对维数 \(k-1\) 成立，则几乎所有 \(\langle T,f,t\rangle\) 都是整数流。
\end{lemma}
\begin{proof}
	由\cref{g3:prop:normal-bounded-flat-precompactness}，先取子列使
	\(\sum_j\mathcal F(T_j-T)<\infty\)。\cref{g3:prop:flat-slicing-estimate}给出
	\[
		\int_{\mathbb R}\sum_j
		\mathcal F\bigl(\langle T_j,f,t\rangle-\langle T,f,t\rangle\bigr)\,\mathrm dt
		\le\operatorname{Lip}(f)\sum_j\mathcal F(T_j-T)<\infty.
	\]
	故几乎每个 \(t\) 上切片平坦收敛，因而弱收敛。另一方面，切片的质量界与边界公式给出
	\[
		\int_{\mathbb R}
		\mathbf N(\langle T_j,f,t\rangle)\,\mathrm dt
		\le\operatorname{Lip}(f)\mathbf N(T_j).
	\]
	Fatou 引理说明几乎每个 \(t\) 都满足
	\(\liminf_j\mathbf N(\langle T_j,f,t\rangle)<\infty\)。
	固定这样的 \(t\)，再选一个正规总量有界的子列。所有切片仍支于 \(K\)，得到第一项结论。若 \(T_j\) 是整数流，\cref{g3:thm:integer-current-slicing}保证其切片为整数流；再用维数 \(k-1\) 的闭包结论即可。
\end{proof}

这一步的两个子列用途不同：第一个使几乎所有参数上的切片同时收敛，第二个在某个已固定的参数处恢复质量与边界质量界。可积性只保证下极限有限，不能把它直接写成整列的逐点一致上界。

\ProseParagraphBreak
设 \(p_\alpha(x)=(x_{i_1},\ldots,x_{i_k})\)，其中 \(i_1<\cdots<i_k\)，是一个坐标 \(k\) 维投影。若前述归纳假设成立，先沿 \(x_{i_1}\) 切一次，再对所得整数 \((k-1)\) 流沿剩余坐标切片。\cref{g3:prop:iterated-current-slicing}与 Fubini 定理便给出
\[
	\langle T,p_\alpha,y\rangle\ \text{在几乎每个 }y\in\mathbb R^k
	\text{ 上是整数零流}.
\]
接下来的两个小节说明，有限个坐标投影中的这些整数点如何决定原流的几何结构。

\subsection{极限的可整流性}
\label{g3:subsec:2205:04}
知道每个切片是有限个点，还需要控制这些点随参数的移动。正规性给出的控制恰好是 BV 型的。先把所需的标量估计写出。

\begin{lemma}\label{g3:lem:slice-bv-pointwise}
	设 \(\nu\) 是 \(\mathbb R^k\) 上的有限正 Radon 测度，\(v\in BV_{\mathrm{loc}}(\mathbb R^k)\)，且 \(|Dv|\le\nu\)。定义
	\[
		\mathcal M\nu(y)=\sup_{r>0}\frac{\nu(B(y,r))}{\omega_kr^k}.
	\]
	则 \(\mathcal M\nu<\infty\) 在 Lebesgue 几乎处处成立，并且在 \(v\) 的任意两个 Lebesgue 点 \(y,z\) 有
	\begin{equation}\label{g3:eq:slice-bv-pointwise}
		|v(y)-v(z)|
		\le C_k|y-z|\bigl(\mathcal M\nu(y)+\mathcal M\nu(z)\bigr).
	\end{equation}
	如果一列函数 \(v_\ell\) 都满足 \(|Dv_\ell|\le\nu\)，则可删去同一个 Lebesgue 零测集，使式\eqref{g3:eq:slice-bv-pointwise}对所有 \(\ell\) 同时成立。
\end{lemma}
\begin{proof}
	对 \(\{\mathcal M\nu>a\}\) 中每一点选取一个满足
	\(\nu(B)>a\mathcal L^k(B)\) 的中心球。在有界部分上应用五倍球覆盖引理，得到互不相交的 \(B_i\)，其五倍球覆盖这些点。因此
	\[
		\mathcal L^k\{\mathcal M\nu>a\}
		\le 5^k\sum_i\mathcal L^k(B_i)
		\le \frac{5^k}{a}\nu(\mathbb R^k).
	\]
	先穷尽有界部分，再令 \(a\to\infty\)，得到几乎处处有限性。

	我们使用球上的 \(L^1\) Poincaré 估计
	\[
		\frac{1}{|B(y,r)|}\int_{B(y,r)}|v-v_{B(y,r)}|\,\mathrm dx
		\le C_kr\,\frac{|Dv|(B(y,2r))}{r^k}.
	\]
	它可以直接从线段积分得到：光滑情形先以
	\(|v(x)-v(w)|\le |x-w|\int_0^1|\nabla v((1-t)x+tw)|\,\mathrm dt\)
	估计球上的双重积分。将 \(t\) 的积分在 \(1/2\) 处分开，后半段交换 \(x,w\)，前半段对 \(x\) 换元时 Jacobian 不超过 \(2^k\)，得到所示估计。局部卷积并令核半径趋零，得到 BV 版本；扩大为二倍球吸收卷积半径，避免球面上的测度原子问题。

	在 Lebesgue 点 \(y\)，把均值按半径 \(r, r/2,r/4,\ldots\) 逐级相减，使用 Poincaré 估计，得到
	\[
		|v(y)-v_{B(y,r)}|
		\le C_k\sum_{j=0}^{\infty}2^{-j}r\,\mathcal M\nu(y)
		\le C_kr\,\mathcal M\nu(y).
	\]
	若 \(d=|y-z|>0\)，则 \(B(z,2d)\subset B(y,4d)\)。先把 \(v(y)\) 与 \(B(y,4d)\) 的均值比较，再把 \(v(z)\) 与 \(B(z,2d)\) 的均值比较；两个均值之差由大球上的 Poincaré 估计控制，至多为 \(C_kd\,\mathcal M\nu(y)\)。三项相加即得式\eqref{g3:eq:slice-bv-pointwise}。最后对可数个函数取 Lebesgue 零测例外集的并。
\end{proof}

\begin{proposition}\label{g3:prop:atomic-slice-graph-cover}
	设 \(1\le k\le n\)，\(T\) 为 \(\mathbb R^n\) 中紧支正规 \(k\) 流。假设对每个坐标 \(k\) 维投影 \(p_\alpha\)，几乎所有
	\(\langle T,p_\alpha,y\rangle\) 都是整数零流。则存在可数个 Lipschitz \(k\) 维图像的并 \(\Gamma\)，使
	\[
		\|T\|(\mathbb R^n\setminus\Gamma)=0,\qquad
		\|T\|\ll\mathcal H^k\llcorner\Gamma.
	\]
\end{proposition}
\begin{proof}
	固定一个坐标投影，重新排列坐标后写作
	\(p(x)=(x_1,\ldots,x_k)\)，记 \(Q_y=\langle T,p,y\rangle\)。
	取紧集 \(K\) 包含 \(\operatorname{spt}T\)。切片质量可积，且几乎处处支于 \(K\cap p^{-1}(y)\)。

	第一步控制切片对标量函数的作用。对
	\(\varphi\in C_c^\infty(\mathbb R^n)\)，满足
	\(\|\varphi\|_\infty\le1\)、\(\|\nabla\varphi\|_\infty\le1\)，令
	\[
		v_\varphi(y)=Q_y(\varphi).
	\]
	由切片积分公式，\(v_\varphi\in L^1(\mathbb R^k)\)。设
	\[
		\beta_i=(-1)^{i-1}\mathrm dx_1\wedge\cdots
		\wedge\widehat{\mathrm dx_i}\wedge\cdots\wedge\mathrm dx_k,
		\qquad
		\mathrm dx_i\wedge\beta_i=\mathrm dx_1\wedge\cdots\wedge\mathrm dx_k.
	\]
	对 \(\eta\in C_c^\infty(\mathbb R^k)\)，展开
	\(\mathrm d(\varphi(\eta\circ p)\beta_i)\)，得到
	\begin{align*}
		\int_{\mathbb R^k}v_\varphi\,\partial_i\eta\,\mathrm dy
		&=T\bigl(\varphi(\partial_i\eta\circ p)\,
		\mathrm dx_1\wedge\cdots\wedge\mathrm dx_k\bigr)\\
		&=\partial T\bigl(\varphi(\eta\circ p)\beta_i\bigr)
		-T\bigl((\eta\circ p)\mathrm d\varphi\wedge\beta_i\bigr).
	\end{align*}
	由于各坐标形式的余质量为一，右侧绝对值至多为
	\[
		\int_{\mathbb R^n}|\eta\circ p|\,\mathrm d(\|T\|+\|\partial T\|).
	\]
	故每个方向导数都由同一个推前测度控制，特别地
	\begin{equation}\label{g3:eq:slice-testing-bv-bound}
		|Dv_\varphi|\le\nu,\qquad
		\nu:=k\,p_\#(\|T\|+\|\partial T\|).
	\end{equation}

	第二步把这个标量估计同时用于全部切片。对支于 \(K\) 的有限有符号测度定义
	\[
		d_{\mathrm{BL}}(A,B)
		=\sup\{|A(\varphi)-B(\varphi)|:
		\|\varphi\|_\infty\le1,\ \operatorname{Lip}(\varphi)\le1\}.
	\]
	这个上确界可以只对可数个光滑函数取。理由是上述函数限制到 \(K\) 后一致有界、等度连续，因而在一致范数中可分；逐个作光滑逼近并以趋于一的因子缩放，可同时保持两个界。测度在 \(K\) 上有限，所以一致逼近不改变上确界。

	对这一可数族应用\cref{g3:lem:slice-bv-pointwise}。删去同一个零测集后，所有 \(y,z\) 均满足
	\begin{equation}\label{g3:eq:slice-configuration-lipschitz}
		d_{\mathrm{BL}}(Q_y,Q_z)
		\le C_k|y-z|\bigl(\mathcal M\nu(y)+\mathcal M\nu(z)\bigr).
	\end{equation}
	因而在每个集合 \(\{\mathcal M\nu\le L\}\) 上，切片配置关于这个距离是 Lipschitz 的。这里的 \(d_{\mathrm{BL}}\) 只是有限测度上的测试函数距离，所需估计已经由式\eqref{g3:eq:slice-testing-bv-bound}证明。

	第三步从配置中取出各个原子的位置。非零整数零流可以写成有限和
	\[
		Q_y=\sum_{\ell=1}^{q}m_\ell\delta_{a_\ell(y)},
		\qquad m_\ell\in\mathbb Z\setminus\{0\},
		\quad a_\ell(y)\ne a_{\ell'}(y)\ (\ell\ne\ell').
	\]
	按以下数据分组：整数 \(L,q\)，有序整数系数 \(m_1,\ldots,m_q\)，以及有理中心、正有理半径的球 \(B_\ell=B(c_\ell,r_\ell)\)，要求二倍球两两不交，每个 \(B_\ell\) 中恰有指定原子，所有原子均在这些球中，并且 \(\mathcal M\nu(y)\le L\)。任意有限个互异原子都可被这样一族球分开，故总共可数个组覆盖了几乎所有非零 \(Q_y\)。

	固定一组以及一个指标 \(\ell\)。取光滑截断 \(\chi_\ell\)，在 \(\overline B_\ell\) 上为一，支于 \(2B_\ell\)。对每个环境坐标 \(s\)，函数
	\[
		\varphi_{\ell,s}(x)=\chi_\ell(x)(x_s-c_{\ell,s})
	\]
	的上确界及 Lipschitz 常数均由所选球控制。因为其支集不遇到其余原子，
	\[
		Q_y(\varphi_{\ell,s})=m_\ell(a_{\ell,s}(y)-c_{\ell,s}).
	\]
	于是对同组的 \(y,z\)，式\eqref{g3:eq:slice-configuration-lipschitz}给出
	\[
		|a_{\ell,s}(y)-a_{\ell,s}(z)|
		\le C\,|y-z|.
	\]
	常数可以依赖于这组数据，但与 \(y,z\) 无关；非零整数保证 \(|m_\ell|\ge1\)。又 \(p(a_\ell(y))=y\)，所以这些点构成参数集合上的 Lipschitz 图。将其余 \(n-k\) 个坐标逐个作 Lipschitz 延拓到整个 \(\mathbb R^k\)，保持前 \(k\) 个坐标等于 \(y\)，得到一个完整的闭 Lipschitz 图像。对全部组和全部原子求并，得到 Borel 集 \(\Gamma_\alpha\)，使几乎每个 \(Q_y\) 都支于 \(\Gamma_\alpha\)。

	最后恢复原流的质量。记 \(\mu_\alpha\) 为 \(T\) 对坐标形式
	\(\mathrm dx_{i_1}\wedge\cdots\wedge\mathrm dx_{i_k}\) 的有符号系数测度。切片积分公式给出
	\begin{equation}\label{g3:eq:coordinate-current-measure-disintegration}
		\mu_\alpha(B)=\int_{\mathbb R^k}Q_y(B)\,\mathrm dy
		\quad(B\subset\mathbb R^n\text{ 为 Borel 集}).
	\end{equation}
	事实上，公式先对连续紧支函数成立；由
	\(\int\mathbf M(Q_y)\,\mathrm dy<\infty\)，单调类与测度唯一性将它推广到 Borel 集。故 \(|\mu_\alpha|\) 集中于 \(\Gamma_\alpha\)。

	若 \(\mathcal H^k(B)=0\)，投影 \(p_\alpha(B)\) 也是 \(k\) 维零测集。对 \(y\notin p_\alpha(B)\)，纤维支集性质给出 \(Q_y(B)=0\)，故 \(|\mu_\alpha|(B)=0\)。有限维向量测度的质量满足
	\(\|T\|\le C_{n,k}\sum_\alpha|\mu_\alpha|\)，反向的每个分量也受 \(\|T\|\) 控制。于是令
	\(\Gamma=\bigcup_\alpha\Gamma_\alpha\)，便得到所述集中性及绝对连续性。
\end{proof}

图覆盖给出了几何集合，还要证明测度中的 \(k\) 向量确实沿着这些图像。这里再次需要边界质量有限。

\begin{lemma}\label{g3:lem:normal-current-tangent-orientation}
	设 \(T\) 为 \(\mathbb R^n\) 中正规 \(k\) 流，\(1\le k\le n\)。若 \(\|T\|\) 集中于可数 Lipschitz \(k\) 图像的并 \(\Gamma\)，且
	\(\|T\|\ll\mathcal H^k\llcorner\Gamma\)，则 \(T\) 具有按 Lipschitz 片定向的实重数表示
	\[
		T(\omega)=\int_E\langle\omega(x),\tau(x)\rangle\theta(x)\,
		\mathrm d\mathcal H^k(x),
	\]
	其中 \(E\subset\Gamma\) 可数 \(k\)-可整流，\(\theta>0\)，而 \(\tau(x)\) 是所在 Lipschitz 片切平面的单位简单取向向量；相交片在交集的几乎处处给出相同切平面。
\end{lemma}
\begin{proof}
	各个完整 Lipschitz 图在有界部分有有限的 \(\mathcal H^k\) 测度，故 \(\mathcal H^k\llcorner\Gamma\) 是 \(\sigma\)-有限测度。Radon--Nikodym 表示与向量测度的极分解给出
	\[
		\|T\|=\theta\,\mathcal H^k\llcorner\Gamma,\qquad
		T=\xi\,\theta\,\mathcal H^k\llcorner\Gamma,\qquad
		|\xi(x)|_{\mathrm m}=1
	\]
	在 \(\theta>0\) 的点几乎处处成立。删去 \(\theta=0\) 的点，记所得集合为 \(E\)。当前只知道 \(\xi\) 是单位质量范数的 \(k\) 向量，还没有它的简单性。

	将图像逐个减去前面图像，得到不交可测片 \(E_i\)，每片包含于一个 Lipschitz 图。对每片使用\cref{g3:thm:rectifiable-tangent-plane,g3:thm:weighted-rectifiable-density}的参数化证明，分别用于 \(\theta\) 和有限个分量 \(\theta\xi_\alpha\)。片外测度在该片的几乎每点具有零 \(k\) 维密度，这是\cref{g3:lem:off-sheet-density-zero}的结论。在交集的密度点，相交图像的切平面相同；可由交集的双 Lipschitz 参数片和式\eqref{g3:eq:tangent-parametric-image}核验。因而对 \(\mathcal H^k\)-几乎每个 \(x\in E_i\)，令 \(P_x\) 为包含该片的图像的切平面，缩放
	\(\eta_{x,r}(z)=(z-x)/r\) 满足
	\begin{equation}\label{g3:eq:current-vector-plane-blowup}
		(\eta_{x,r})_\#T\rightharpoonup S_x,\qquad
		S_x(\omega)=\theta(x)\int_{P_x}\langle\omega(z),\xi(x)\rangle\,
		\mathrm d\mathcal H^k(z).
	\end{equation}
	这些应用逐片进行，故不要求整个并集 \(\Gamma\) 在有界区域内有有限的 \(\mathcal H^k\) 测度。

	令 \(\lambda=\|\partial T\|\)。有限测度的覆盖估计给出
	\begin{equation}\label{g3:eq:boundary-mass-vanishing-at-k-points}
		\lim_{r\downarrow0}\frac{\lambda(B(x,r))}{r^{k-1}}=0
		\quad\text{在 }\mathcal H^k\text{-几乎每个 }x\text{ 成立}.
	\end{equation}
	说明如下：当 \(k=1\) 时，非零极限只能出现在 \(\lambda\) 的原子上，原子至多可数。当 \(k>1\) 时，对
	\(\{\limsup_{r\downarrow0}\lambda(B(x,r))/r^{k-1}>1/a\}\)
	的有界部分，在任意小尺度选满足该下界的球并作五倍球覆盖，得到其 \(\mathcal H^{k-1}\) 测度至多为 \(C_ka\,\lambda(\mathbb R^n)\)。因此坏集合是 \(\mathcal H^{k-1}\)-有限集的可数并，其 \(\mathcal H^k\) 测度为零。

	取同时满足式\eqref{g3:eq:current-vector-plane-blowup}与式\eqref{g3:eq:boundary-mass-vanishing-at-k-points}的 \(x\)。对支于 \(B_R\) 的测试 \((k-1)\) 形式 \(\zeta\)，有
	\[
		|\partial((\eta_{x,r})_\#T)(\zeta)|
		\le \|\zeta\|_\infty r^{1-k}\lambda(B(x,Rr))
		\longrightarrow0.
	\]
	所以 \(\partial S_x=0\)。

	在正交坐标中把 \(P_x\) 写成 \(\mathbb R^k\times\{0\}\)。若 \(\ell>k\)，对任意不含 \(\ell\) 的 \((k-1)\) 元多重指标 \(J\)，取
	\(\zeta(z)=z_\ell\psi(z)\,\mathrm dz_J\)，其中 \(\psi\in C_c^\infty(\mathbb R^n)\)。在 \(P_x\) 上有
	\(\mathrm d\zeta=\psi\,\mathrm dz_\ell\wedge\mathrm dz_J\)，于是
	\[
		0=\partial S_x(\zeta)
		=\theta(x)\langle\mathrm dz_\ell\wedge\mathrm dz_J,\xi(x)\rangle
		\int_{P_x}\psi\,\mathrm d\mathcal H^k.
	\]
	选择平面积分非零的 \(\psi\)，便知 \(\xi(x)\) 中每个含法向指标的系数为零。因此
	\(\xi(x)\in\Lambda_kP_x\)，而这个空间是一维的。单位质量范数遂给出
	\(\xi(x)=\pm e_1\wedge\cdots\wedge e_k\)。这就是所需的简单切向取向 \(\tau(x)\)，并完成按片的定向表示。
\end{proof}

\subsection{整数重数的保持}
\label{g3:subsec:2205:05}
现在已有定向片状表示。投影的 Jacobian 把每片上的面积密度换算为纤维中的原子系数，因而可以逐点读出重数。

\begin{theorem}[原子切片判别]\label{g3:thm:atomic-slice-rectifiability}
	设 \(1\le k\le n\)，\(T\) 为 \(\mathbb R^n\) 中紧支正规 \(k\) 流。如果对每个坐标 \(k\) 维投影 \(p_\alpha\)，几乎所有零维切片
	\(\langle T,p_\alpha,y\rangle\) 都是整数零流，则 \(T\) 是整数可整流流。
\end{theorem}
\begin{proof}
	\cref{g3:prop:atomic-slice-graph-cover,g3:lem:normal-current-tangent-orientation}已给出
	\[
		T=\tau\theta\,\mathcal H^k\llcorner E,\qquad \theta>0,
	\]
	其中 \(E\) 是可数 Lipschitz 片的并，\(\tau\) 是片上的单位简单切向向量。固定坐标投影 \(p=p_\alpha\)，令
	\[
		a_\alpha(x)=\langle\mathrm dx_{i_1}\wedge\cdots
		\wedge\mathrm dx_{i_k},\tau(x)\rangle,\qquad
		J_\alpha(x)=|a_\alpha(x)|.
	\]
	对可数个不交参数片使用面积公式，得到
	\begin{align*}
		T(\varphi\,\mathrm dx_\alpha)
		&=\int_E\varphi(x)\theta(x)a_\alpha(x)\,\mathrm d\mathcal H^k(x)\\
		&=\int_{\mathbb R^k}
		\sum_{\substack{x\in E\cap p^{-1}(y)\\J_\alpha(x)>0}}
		\varphi(x)\theta(x)\operatorname{sgn}(a_\alpha(x))\,\mathrm dy.
	\end{align*}
	其中 \(J_\alpha=0\) 的部分不贡献第一行积分。取绝对值并再用面积公式，知右侧各原子系数的绝对值之和对 \(y\) 可积。将测试函数换成 \((g\circ p_\alpha)\varphi\)，其中 \(g\in C_c^\infty(\mathbb R^k)\)，再与式\eqref{g3:eq:coordinate-current-measure-disintegration}比较。对每个固定 \(\varphi\)，两个纤维配对作为 \(y\) 的可积函数在所有 \(g\) 上积分相同，因而几乎处处相等。对可数稠密的 \(\varphi\) 族删去同一个零测集，有限测度的唯一性便给出
	\begin{equation}\label{g3:eq:slice-reads-integer-multiplicity}
		\langle T,p_\alpha,y\rangle
		=\sum_{\substack{x\in E\cap p_\alpha^{-1}(y)\\J_\alpha(x)>0}}
		\theta(x)\operatorname{sgn}(a_\alpha(x))\,\delta_x
		\quad\text{对几乎所有 }y.
	\end{equation}
	不同的 \(x\) 是不同原子，故不能彼此抵消。左侧为整数零流，因而右侧每个非零系数都是整数。

	记 \(B=\{x\in E:\theta(x)\notin\mathbb N\}\)。由刚才的结论与面积公式，
	\[
		\int_BJ_\alpha\,\mathrm d\mathcal H^k=0
		\qquad\text{对每个 }\alpha.
	\]
	单位简单向量的坐标系数满足
	\(\sum_\alpha a_\alpha^2=1\)，所以有限个 \(J_\alpha\) 不会同时为零。于是 \(\mathcal H^k(B)=0\)，即重数为正整数。特别地，
	\[
		\mathcal H^k(E\cap U)
		\le\int_{E\cap U}\theta\,\mathrm d\mathcal H^k
		=\|T\|(U)
	\]
	对每个开集 \(U\) 成立。此时 \(E\) 的 Hausdorff 测度局部有限，前章的近似切平面定理直接适用于整个 \(E\)，并与各片的切平面一致。因此得到整数可整流流的完整表示。
\end{proof}

这也完成了维数归纳所缺的一步：零维切片的整数信息先产生图覆盖，再决定图上的取向与重数。

\begin{theorem}[整数流闭包]\label{g3:thm:integer-current-closure}
	设 \(K\subset\mathbb R^n\) 紧，\(0\le k\le n\)，\(T_j\in\mathcal I_k(\mathbb R^n)\)，并且
	\[
		\operatorname{spt}T_j\subset K,\qquad
		\sup_j\bigl(\mathbf M(T_j)+\mathbf M(\partial T_j)\bigr)<\infty.
	\]
	如果 \(T_j\rightharpoonup T\)，则 \(T\in\mathcal I_k(\mathbb R^n)\)。
\end{theorem}
\begin{proof}
	按 \(k\) 归纳。\(k=0\) 正是\cref{g3:lem:integer-zero-current-closure}。设结论已对维数 \(k-1\) 成立。质量下半连续性与测试形式的边界公式先保证 \(T\) 正规、支于 \(K\)。

	对每个坐标 \(k\) 维投影 \(p_\alpha\)，先沿它的第一个坐标切片。\cref{g3:lem:slice-closure-subsequence}和归纳假设说明，几乎所有这些切片是整数 \((k-1)\) 流。对它们再沿其余坐标作迭代切片，得到几乎所有 \(\langle T,p_\alpha,y\rangle\) 都是整数零流。这里的零测例外集由切片公式和可数测试函数选取，Fubini 定理可把逐次的“几乎处处”合并为 \(\mathbb R^k\) 中的“几乎处处”。坐标投影只有有限个，故可以分别完成此论证。\cref{g3:thm:atomic-slice-rectifiability}于是说明 \(T\) 整数可整流。

	还须处理边界。各 \(\partial T_j\) 是整数 \((k-1)\) 流，支于 \(K\)，质量一致有界，而二次边界均为零。又
	\(\partial T_j\rightharpoonup\partial T\)。再次使用维数 \(k-1\) 的归纳假设，得到 \(\partial T\) 也是整数流，特别地是整数可整流流。因此 \(T\) 满足整数流定义中的两项要求，归纳完成。
\end{proof}

\subsection{边界相容及下半连续}
\label{g3:subsec:2205:06}
闭包定理保留了竞争者的种类。对于极小化，还需要极限保留指定边界，并使目标质量不增加。

\begin{corollary}\label{g3:cor:integral-limit-boundary-lsc}
	在\cref{g3:thm:integer-current-closure}的条件下，如果
	\(\partial T_j=B_j\rightharpoonup B\)，则 \(\partial T=B\)。对任意非负连续紧支函数 \(\phi\)，有
	\[
		\int\phi\,\mathrm d\|T\|
		\le\liminf_j\int\phi\,\mathrm d\|T_j\|,\qquad
		\int\phi\,\mathrm d\|\partial T\|
		\le\liminf_j\int\phi\,\mathrm d\|\partial T_j\|.
	\]
	特别地，固定整数边界 \(\partial T_j=B\) 的条件在这种极限下保持。
\end{corollary}
\begin{proof}
	对每个紧支光滑 \((k-1)\) 形式 \(\eta\)，有
	\[
		\partial T(\eta)
		=T(\mathrm d\eta)
		=\lim_jT_j(\mathrm d\eta)
		=\lim_j B_j(\eta)=B(\eta).
	\]
	这说明边界相容不需要额外选取子列。对质量，开集版本的下半连续性已在式\eqref{g3:eq:normal-current-limit-mass}证明。由于 \(\{\phi>t\}\) 为开集，层蛋糕表示与 Fatou 引理给出
	\[
		\int\phi\,\mathrm d\|T\|
		=\int_0^\infty\|T\|(\{\phi>t\})\,\mathrm dt
		\le\liminf_j\int_0^\infty\|T_j\|(\{\phi>t\})\,\mathrm dt.
	\]
	边界质量同理。取在共同紧支集邻域等于一的 \(\phi\)，便恢复总质量与边界总质量的下半连续性。
\end{proof}

固定边界与质量下半连续性具有不同来源：前者来自 \(\partial T(\eta)=T(\mathrm d\eta)\)，后者来自单位测试形式的上确界。整数闭包则用到了切片、BV 估计和切向结构；三者合在一起，才得到可用于面积极小化的极限。

\subsection{整数流的整体与局部紧性}
\label{g3:subsec:2205:07}
把弱紧性、平坦控制与整数闭包合并，就得到本章的主要紧性定理。

\begin{theorem}[Federer--Fleming 紧性定理]\label{g3:thm:federer-fleming-compactness}
	设 \(0\le k\le n\)，\(K\subset\mathbb R^n\) 紧。若整数流 \(T_j\in\mathcal I_k(\mathbb R^n)\) 满足
	\[
		\operatorname{spt}T_j\subset K,\qquad
		\sup_j\bigl(\mathbf M(T_j)+\mathbf M(\partial T_j)\bigr)<\infty,
	\]
	则存在子列与 \(T\in\mathcal I_k(\mathbb R^n)\)，使
	\[
		\operatorname{spt}T\subset K,\qquad
		\mathcal F(T_j-T)\longrightarrow0,\qquad
		T_j\rightharpoonup T,\quad
		\partial T_j\rightharpoonup\partial T.
	\]
	极限的质量及边界质量满足局部与整体的下半连续性。
\end{theorem}
\begin{proof}
	\cref{g3:thm:normal-current-weak-compactness}先给出支于 \(K\) 的正规弱极限。\cref{g3:thm:integer-current-closure}保证它属于 \(\mathcal I_k\)。\cref{g3:prop:normal-bounded-flat-precompactness}把同一子列的弱收敛加强为平坦收敛；边界相容及下半连续性由\cref{g3:cor:integral-limit-boundary-lsc}给出。
\end{proof}

\begin{theorem}\label{g3:thm:local-integral-current-compactness}
	设 \(\Omega\subset\mathbb R^n\) 开，\(T_j\) 是 \(\Omega\) 中局部整数 \(k\) 流，即它们及其边界在每个相对紧开集上都具有局部有限质量的整数可整流表示。若对每个 \(V\Subset\Omega\) 都有
	\[
		\sup_j\bigl(\|T_j\|(V)+\|\partial T_j\|(V)\bigr)<\infty,
	\]
	则可取子列，使 \(T_j\rightharpoonup T\)，其中 \(T\) 为局部整数流。对每个 \(V\Subset\Omega\)，取任意 \(\eta\in C_c^\infty(\Omega)\) 在 \(V\) 的邻域等于一，则还有
	\[
		\mathcal F\bigl(\eta T_j-\eta T\bigr)\longrightarrow0.
	\]
	这里把紧支流补零到 \(\mathbb R^n\) 后计算平坦范数。
\end{theorem}
\begin{proof}
	先用可数个相对紧开集穷尽 \(\Omega\)，在每个紧支测试形式空间上重复弱紧性的可数稠密族论证，再取对角子列。得到局部正规流 \(T\)，且 \(T_j\rightharpoonup T\)，局部质量与边界质量下半连续。

	为证明整数性，取 \(\overline{B_R(x)}\subset\Omega\)。在一个略大的紧集上，各正测度
	\(\mu_j=\|T_j\|+\|\partial T_j\|\) 的质量一致有界。取子列使它们在该紧集上弱星收敛于有限正测度 \(\mu\)。由局部下半连续性，\(\|T\|+\|\partial T\|\le\mu\) 在球内成立。除至多可数个半径外，\(\mu(\partial B_r(x))=0\)。

	对这样的 \(r<R\)，用光滑径向函数从两侧逼近 \(\mathbf1_{B_r(x)}\)，得到
	\[
		T_j\llcorner B_r(x)\rightharpoonup T\llcorner B_r(x),\qquad
		(\partial T_j)\llcorner B_r(x)\rightharpoonup
		(\partial T)\llcorner B_r(x).
	\]
	具体地，逼近误差支于厚度为 \(\delta\) 的球面邻域，其绝对值由测试形式的一致范数乘该邻域的 \(\mu_j\) 质量控制。先令 \(j\to\infty\)，再令 \(\delta\downarrow0\)，误差因 \(\mu(\partial B_r(x))=0\) 而消失。

	如果 \(k=0\)，这些限制流已经是共同紧支、质量有界的整数零流，直接用\cref{g3:lem:integer-zero-current-closure}即可。以下设 \(k\ge1\)。对距离函数 \(d_x(z)=|z-x|\) 使用切片质量估计和 Fatou 引理，有
	\[
		\int_{R/2}^{R}\liminf_j
		\mathbf M(\langle T_j,d_x,r\rangle)\,\mathrm dr
		\le\sup_j\|T_j\|(B_R(x))<\infty.
	\]
	因而可以选 \(r\in(R/2,R)\)，同时使上述限制收敛成立、每个 \(T_j\llcorner B_r(x)\) 是整数流，并使切片质量的下极限有限。再取一个可依赖于 \(r\) 的子列，使该切片质量一致有界。由
	\[
		\partial(T_j\llcorner B_r(x))
		=(\partial T_j)\llcorner B_r(x)+\langle T_j,d_x,r\rangle
	\]
	可知这些限制流的正规总量一致有界，且共同支于 \(\overline{B_r(x)}\)。固定紧支的整数闭包定理于是给出 \(T\llcorner B_r(x)\) 为整数流。在 \(B_r(x)\) 内，它及其边界与 \(T,\partial T\) 分别一致。用可数个这样的内球覆盖 \(\Omega\)，便得到 \(T\) 的局部整数性；不同球中使用的子列只用于识别同一个已经得到的弱极限。

	最后，\(\eta T_j\) 一般为实系数流，但仍是共同紧支的正规流。由乘积公式，其正规总量受
	\[
		\|\eta\|_\infty
		\bigl(\|T_j\|(\operatorname{spt}\eta)+\|\partial T_j\|(\operatorname{spt}\eta)\bigr)
		+\|\mathrm d\eta\|_\infty\|T_j\|(\operatorname{spt}\eta)
	\]
	一致控制。又 \(\eta T_j\rightharpoonup\eta T\)，故\cref{g3:prop:normal-bounded-flat-precompactness}给出所述平坦收敛。
\end{proof}

局部结论容许质量逃向无穷远。例如在 \(\mathbb R^2\) 中，把逆时针单位圆流 \(C\) 平移到圆心 \((j,0)\)，得到 \(C_j\)。始终有
\[
	\mathbf M(C_j)=2\pi,\qquad \partial C_j=0,
\]
但每个紧支测试形式最终都遇不到 \(C_j\)，故 \(C_j\rightharpoonup0\)。它们却不在全空间平坦范数中趋于零：取在单位圆附近等于
\(\tfrac12(-x_2\,\mathrm dx_1+x_1\,\mathrm dx_2)\) 的紧支光滑形式 \(\omega\)，则 \(C(\omega)=\pi\)；平移得到的 \(\omega_j\) 具有与 \(j\) 无关的余质量和外微分余质量上界，所以
\[
	\mathcal F(C_j)\ge
	\frac{\pi}{\max\{\|\omega\|_\infty,\|\mathrm d\omega\|_\infty\}}>0.
\]
共同紧支条件排除了这种逃逸。下一章在选择极小化序列时，还需要给出把竞争者留在一个共同有界区域中的几何理由。
% 纲要条目排为编号小节标题，并进入目录与书签；说明段仍保留于纲要注释中。


\begin{chaptersummary}
	一般流是测试形式上的连续线性泛函，有限质量流由向量 Radon 测度表示。余质量与其对偶范数使光滑定向曲面的质量等于面积。边界定义自动给出 \(\partial^2=0\)；正规性同时控制流与边界的质量，使光滑推前、乘积和同伦公式能够延伸到 Lipschitz 映射。

	\ProseParagraphBreak
	可整流流还要求向量沿可整流片的切平面，整数流则要求自身及其边界都具有整数重数的可整流表示。有限周长集合给出顶维模型，整数值 BV 函数通过层集分解产生更多例子。取向改变可以形成内部边界；弱极限中的抵消则可能使质量严格下降。

	\ProseParagraphBreak
	切片是限制流后新增的边界。其边界公式每降一维带一个负号，质量积分受切割函数的 Lipschitz 常数控制。锥把闭流连接到一个顶点，产生高一维填充；零维时还须消除总系数。平坦范数允许留下小的本体误差，并用低质量流填充其余部分，因此比质量范数更适合描述几何移动。

	\ProseParagraphBreak
	网格变形先选择质量稀疏的径向中心，以可积的距离奇性控制推前和同伦，再用面内 BV 中位数修正得到常数重数的多面体流。统一的尺度估计把正规流有界族逼近为有限维有界集合。闭包证明随后利用整数零维切片：共同的 BV 估计产生可数 Lipschitz 图覆盖，边界质量使放大极限向量沿这些图的切平面，面积公式确定整数重数。结合弱紧性、边界相容与质量下半连续性，得到整数流的平坦紧性。局部版本仍容许质量离开每个固定紧集；下一章的 Plateau 问题将另行处理竞争者的共同支集。
\end{chaptersummary}

\BookExercises

\begin{exercise}\label{g3:ex:weighted-triangle-current}
	设 \(A,B,C\in\mathbb R^2\) 为不共线三点，\(a,b,c\in\mathbb R\)，令
	\[
		T=a\lBrack A,B\rBrack+b\lBrack B,C\rBrack+c\lBrack C,A\rBrack.
	\]
	求 \(\partial T\)、\(\mathbf M(T)\) 和 \(\mathbf M(\partial T)\)。证明 \(\partial T=0\) 当且仅当 \(a=b=c\)。再确定 \(T\) 为整数流的充要条件。
\end{exercise}
\begin{solution}
	逐条使用线段的边界公式，得到
	\[
		\partial T=(c-a)\delta_A+(a-b)\delta_B+(b-c)\delta_C.
	\]
	三条线段除端点外不交，端点的一维测度为零，故
	\[
	\begin{aligned}
		\mathbf M(T)&=|a|\,|B-A|+|b|\,|C-B|+|c|\,|A-C|,\\
		\mathbf M(\partial T)&=|c-a|+|a-b|+|b-c|.
	\end{aligned}
	\]
	三个点流线性无关，因为可分别在各点附近取支集互不相交的测试函数。因此边界为零当且仅当三个差同时为零，即 \(a=b=c\)。

	若 \(a,b,c\) 为整数，流与边界显然都是整数重数可整流流。反之，在任一边内部取不遇其他两边的小球；此处的流相对于该边的单位取向具有常数密度，依次为 \(a,b,c\)。整数可整流表示的唯一性迫使它们为整数。所以 \(T\in\mathcal I_1\) 当且仅当 \(a,b,c\in\mathbb Z\)。仅有边界系数为整数不够：例如 \(a=b=c=1/2\) 的边界为零，但流不是整数流。
\end{solution}

\begin{exercise}\label{g3:ex:pushforward-degree-and-fold}
	标准单位圆 \(S^1\subset\mathbb R^2\) 取逆时针方向。对整数 \(m\)，令 \(f_m(e^{it})=e^{imt}\)，在圆的邻域内作光滑延拓。求 \(f_{m\#}\lBrack S^1\rBrack\) 及其质量。再令 \(g:\mathbb R\to\mathbb R\)，\(g(x)=|x|\)，求 \(g_\#\lBrack-1,1\rBrack\)，解释面积公式中的绝对重数与流的带符号重数的区别。
\end{exercise}
\begin{solution}
	将测试一形式沿圆写成 \(a(t)\,\mathrm dt\)，其中 \(a\) 为 \(2\pi\) 周期函数。参数计算给出
	\[
		f_{m\#}\lBrack S^1\rBrack(\omega)
		=\int_0^{2\pi}m\,a(mt)\,\mathrm dt
		=m\int_0^{2\pi}a(t)\,\mathrm dt.
	\]
	正负整数都可按换元及周期性检验，\(m=0\) 时拉回一形式为零。故推前是 \(m\lBrack S^1\rBrack\)，质量为 \(2\pi|m|\)。当 \(m\ne0\)，几乎每个像点有 \(|m|\) 个原像，取向贡献的符号全部相同。

	对目标直线上的测试形式 \(a(y)\,\mathrm dy\)，
	\[
	\begin{aligned}
		g_\#\lBrack-1,1\rBrack(a\,\mathrm dy)
		&=-\int_{-1}^0a(-x)\,\mathrm dx+\int_0^1a(x)\,\mathrm dx=0.
	\end{aligned}
	\]
	所以推前流为零。相应的绝对 Jacobian 积分却是 \(\int_{-1}^1|g'|=2\)：\((0,1)\) 中每个点有两个原像，其方向相反。面积公式计两份长度，流的重数则是 \(-1+1=0\)。两端点原本的边界 \(\delta_1-\delta_{-1}\) 也都映到一点而抵消，与边界相容一致。
\end{solution}

\begin{exercise}\label{g3:ex:integer-bv-concentric-current}
	设 \(n\ge2\)，\(0<r<R\)，并令
	\[
		u=\chi_{B_R}-2\chi_{B_r},\qquad
		T=u\lBrack\mathbb R^n\rBrack.
	\]
	证明 \(T\) 是整数流，计算其质量、边界及边界质量。若以相同支集上的正重数函数 \(|u|\) 代替 \(u\)，这些量怎样变化？
\end{exercise}
\begin{solution}
	由有限周长集合的流表示，
	\[
		T=\lBrack B_R\rBrack-2\lBrack B_r\rBrack,
		\qquad
		\partial T=\lBrack\partial B_R\rBrack-2\lBrack\partial B_r\rBrack.
	\]
	球面均取各自球的外法向优先取向。两个边界互不相交，系数为整数，所以流与边界均为整数可整流流。函数 \(u\) 在小球内为 \(-1\)，在环域中为 \(1\)，在大球外为零。因此
	\[
		\mathbf M(T)=\omega_nR^n,
		\qquad \mathbf M(\partial T)
		=n\omega_nR^{n-1}+2n\omega_nr^{n-1}.
	\]
	内球面系数为 \(-2\)，记录从内侧值 \(-1\) 到外侧值 \(1\) 的取向改变；用 \(D\chi_B=-\nu_B\mathcal H^{n-1}\llcorner\partial B\) 也能核对同一符号。

	将密度换为 \(|u|=\chi_{B_R}\)，流变为 \(\lBrack B_R\rBrack\)，质量仍为 \(\omega_nR^n\)，边界却只剩 \(\lBrack\partial B_R\rBrack\)，质量为 \(n\omega_nR^{n-1}\)。可见支集和质量不能确定内部取向跳变，也不能确定边界。
\end{solution}

\begin{exercise}\label{g3:ex:mobius-cut-boundary}
	固定 \(a>0\)，取一个标准 Möbius 带的光滑参数化 \(F:[0,2\pi]\times[-a,a]\to M\)，满足 \(F(2\pi,t)=F(0,-t)\)，并在矩形内部为嵌入。矩形按 \(\mathrm ds\wedge\mathrm dt\) 定向，令 \(T=F_\#\lBrack[0,2\pi]\times[-a,a]\rBrack\)，再令横截线 \(\sigma(t)=F(0,t)\) 按 \(t\) 增大定向。计算 \(\partial T\) 中来自两条竖边的项，说明它为何不消失。
\end{exercise}
\begin{solution}
	矩形在 \(s=2\pi\) 的边界方向为 \(+\partial_t\)，在 \(s=0\) 的方向为 \(-\partial_t\)。因此两条竖边推前的总贡献为
	\[
		F(2\pi,\cdot)_\#\lBrack-a,a\rBrack
		-F(0,\cdot)_\#\lBrack-a,a\rBrack.
	\]
	因 \(F(2\pi,t)=\sigma(-t)\)，第一项为 \(-\sigma_\#\lBrack-a,a\rBrack\)，第二项也为这一流的负号。故总贡献为
	\[
		-2\sigma_\#\lBrack-a,a\rBrack.
	\]
	其余两条水平边映到 Möbius 带的几何边界；竖边识别后的横截线位于带的内部，却在整数流边界中保留了两份贡献。若粘合没有反转 \(t\)，两条竖边的符号才会相消。这个计算把非定向障碍具体表现为内部边界项。
\end{solution}

\begin{exercise}\label{g3:ex:graph-current-slice-signs}
	令 \(F(s,t)=(s,t,st)\)，\(0\le s,t\le1\)，并令
	\[
		T=F_\#\lBrack[0,1]^2\rBrack,
	\]
	方形取 \(\mathrm ds\wedge\mathrm dt\) 的取向。计算 \(0<a,b<1\) 时的 \(\langle T,x,a\rangle\)、\(\langle T,(x,y),(a,b)\rangle\) 及 \(\langle T,(y,x),(b,a)\rangle\)。直接核对第一种切片的边界符号，并计算 \(\int_{\mathbb R}\mathbf M(\langle T,x,a\rangle)\,\mathrm da\)。将该积分与图面的面积元及切向梯度 \(|\nabla_Mx|\) 比较。
\end{exercise}
\begin{solution}
	图面的参数取向由
	\[
		F_s=(1,0,t),\qquad F_t=(0,1,s),\qquad
		F_s\wedge F_t
	\]
	给出。因为 \(\mathrm dx(F_s)=1\)、\(\mathrm dx(F_t)=0\)，沿 \(x=a\) 的切片是从 \((a,0,0)\) 到 \((a,1,a)\) 的直线段，取 \(t\) 增大的方向。因此
	\[
		\langle T,x,a\rangle=(t\mapsto(a,t,at))_\#\lBrack[0,1]\rBrack,
		\qquad
		\partial\langle T,x,a\rangle=\delta_{(a,1,a)}-\delta_{(a,0,0)}.
	\]
	方形的下边取 \(+s\) 方向，上边取 \(-s\) 方向，故其图像的边界被 \(x=a\) 切割时，下边贡献正点，上边贡献负点；另外两条边的 \(x\) 值固定为零或一。于是
	\[
		\langle\partial T,x,a\rangle
		=\delta_{(a,0,0)}-\delta_{(a,1,a)}
		=-\partial\langle T,x,a\rangle.
	\]
	再沿 \(y=b\) 切上述有向线段，得到
	\[
		\langle T,(x,y),(a,b)\rangle=\delta_{(a,b,ab)}.
	\]
	交换切割顺序时，\(\mathrm dy\wedge\mathrm dx=-\mathrm dx\wedge\mathrm dy\)，所以
	\[
		\langle T,(y,x),(b,a)\rangle=-\delta_{(a,b,ab)}.
	\]
	也可以直接看出：先沿 \(y=b\) 切片时，水平参数曲线取 \(-s\) 方向，再按 \(x=a\) 切割便产生负点。

	当 \(0<a<1\) 时，第一种线段切片的长度为 \(\sqrt{1+a^2}\)，其余几乎所有切片为空，故
	\[
		\int_{\mathbb R}\mathbf M(\langle T,x,a\rangle)\,\mathrm da
		=\int_0^1\sqrt{1+a^2}\,\mathrm da
		=\frac12\bigl(\sqrt2+\log(1+\sqrt2)\bigr).
	\]
	参数度量为
	\[
		g=\begin{pmatrix}1+t^2&st\\st&1+s^2\end{pmatrix},\qquad
		\det g=1+s^2+t^2.
	\]
	因此面积元为 \(\sqrt{1+s^2+t^2}\,\mathrm ds\,\mathrm dt\)，而
	\[
		|\nabla_Mx|^2=g^{11}=\frac{1+s^2}{1+s^2+t^2}.
	\]
	两者相乘后得到 \(\sqrt{1+s^2}\,\mathrm ds\,\mathrm dt\)，积分恰为刚算出的切片质量积分。这说明切片计算的量是带切向 Jacobian 的面积，而非未经加权的全部面积。
\end{solution}

\begin{exercise}\label{g3:ex:flat-distance-two-points}
	设 \(a,b\in\mathbb R^n\)，证明
	\[
		\mathcal F(\delta_b-\delta_a)=\min\{|b-a|,2\},\qquad
		\mathcal F(\delta_a)=1.
	\]
	解释第一个公式中的两种上界分别对应什么平坦分解，并说明为什么单个正点不能由紧支一维流填充。
\end{exercise}
\begin{solution}
	若 \(a=b\)，第一式两边都为零。以下设 \(d=|b-a|>0\)。取 \(R=\delta_b-\delta_a\)、\(S=0\)，得到上界二；取 \(R=0\)、\(S\) 为从 \(a\) 到 \(b\) 的有向线段，得到上界 \(d\)。前者保留两个端点作为本体误差，后者完全用线段填充。

	为证下界，令 \(e=(b-a)/d\)、\(m=(a+b)/2\)，取有界的 \(1\)-Lipschitz 函数
	\[
		\psi(x)=\max\{-1,\min\{1,e\cdot(x-m)\}\}.
	\]
	则 \(|\psi|\le1\)，且 \(\psi(b)-\psi(a)=\min\{d,2\}\)。将它乘一个在两点附近等于一、梯度至多 \(\varepsilon\) 的大尺度紧支截断，所得函数的 Lipschitz 常数不超过 \(1+\varepsilon\)。再作足够小半径的非负卷积，并除以 \(1+\varepsilon\)，便得到满足 \(\|\varphi\|_\infty\le1\)、\(\|\mathrm d\varphi\|_\infty\le1\) 的光滑紧支测试函数。令卷积半径和 \(\varepsilon\) 趋零，其两点值之差趋于 \(\min\{d,2\}\)。由平坦对偶公式得到所需下界。

	对单个正点，取 \(R=\delta_a,S=0\) 得上界一。另取在 \(a\) 附近等于一、取值在 \([0,1]\) 中且梯度至多一的光滑紧支函数，则对偶公式给出下界一。

	若紧支一维流 \(S\) 满足 \(\partial S=\delta_a\)，取在 \(\operatorname{spt}S\cup\{a\}\) 邻域等于一的紧支测试函数 \(\eta\)，便有
	\[
		1=\delta_a(\eta)=\partial S(\eta)=S(\mathrm d\eta)=0,
	\]
	矛盾。因此零维边界的总系数为零是紧支填充的必要条件。
\end{solution}

\begin{exercise}\label{g3:ex:annular-current-cancellation}
	在 \(\mathbb R^2\) 中取标准取向。对 \(0<\varepsilon<1\)，令
	\[
		A_\varepsilon=\{x:1<|x|<1+\varepsilon\},\qquad
		S_\varepsilon=\lBrack A_\varepsilon\rBrack,\qquad
		T_\varepsilon=\partial S_\varepsilon.
	\]
	求 \(T_\varepsilon\) 两个圆周分支的取向、\(\mathbf M(T_\varepsilon)\) 及 \(\mathbf M(\partial T_\varepsilon)\)。证明 \(T_\varepsilon\) 平坦收敛于零，并证明其质量测度弱星收敛于 \(2\mathcal H^1\llcorner\partial B_1\)。说明这与流的弱极限为何不同。
\end{exercise}
\begin{solution}
	环域的外圆周取逆时针方向，内圆周取顺时针方向。因此若 \(C_r\) 表示半径 \(r\) 的逆时针有向圆周，则
	\[
		T_\varepsilon=C_{1+\varepsilon}-C_1,
		\qquad \mathbf M(T_\varepsilon)=2\pi(1+\varepsilon)+2\pi,
		\qquad \partial T_\varepsilon=0.
	\]
	由于 \(T_\varepsilon=\partial S_\varepsilon\)，平坦范数满足
	\[
		\mathcal F(T_\varepsilon)\le\mathbf M(S_\varepsilon)
		=\pi\bigl((1+\varepsilon)^2-1\bigr)\longrightarrow0.
	\]
	所以流 \(T_\varepsilon\) 弱收敛于零，但它的质量趋于 \(4\pi\)，并不趋于零流的质量。

	对任意连续紧支函数 \(\varphi\)，极坐标参数给出
	\[
	\begin{aligned}
		\int\varphi\,\mathrm d\|T_\varepsilon\|
		&=\int_0^{2\pi}\varphi((1+\varepsilon)(\cos t,\sin t))(1+\varepsilon)\,\mathrm dt\\
		&\quad+\int_0^{2\pi}\varphi(\cos t,\sin t)\,\mathrm dt\\
		&\longrightarrow2\int_{\partial B_1}\varphi\,\mathrm d\mathcal H^1.
	\end{aligned}
	\]
	两种极限记录的信息不同：流的配对保留圆周取向，重合后两个相反的取向抵消；质量测度先取绝对值，所以两份长度相加。这个例子同时满足共同紧支与统一正规质量界，说明这些紧性条件仍不保证质量在极限中保持。
\end{solution}

\begin{exercise}\label{g3:ex:cubical-staircase-flat-mass}
	在 \(\mathbb R^2\) 中，令 \(T\) 为从 \((0,0)\) 指向 \((1,1)\) 的对角线段积分流。对正整数 \(m\)，令 \(h=1/m\)，以折线路径
	\[
	(ih,ih)\longrightarrow((i+1)h,ih)
	\longrightarrow((i+1)h,(i+1)h),\qquad 0\le i<m,
	\]
	定义网格积分流 \(P_m\)。计算 \(\mathbf M(T)\)、\(\mathbf M(P_m)\) 及 \(\mathbf M(P_m-T)\)，并构造整数 \(2\) 流 \(S_m\)，证明
	\[
	\partial S_m=P_m-T,\qquad \mathcal F(P_m-T)\le\frac1{2m}.
	\]
	进一步证明，在任意坐标立方网格上，凡有限多面体 \(1\) 流 \(P\) 满足
	\(\partial P=\delta_{(1,1)}-\delta_{(0,0)}\)，都有 \(\mathbf M(P)\ge2\)。据此说明，固定坐标网格的平坦逼近为何不必保持质量。
\end{exercise}
\begin{solution}
	对角线的长度为 \(\sqrt2\)，而折线有 \(m\) 个水平段和 \(m\) 个竖直段，每段长 \(h\)，因此
	\[
	\mathbf M(T)=\sqrt2,\qquad \mathbf M(P_m)=2.
	\]
	两条路径的交集只有有限个顶点，具有零 \(\mathcal H^1\) 测度；相应质量测度互相奇异，所以
	\(\mathbf M(P_m-T)=2+\sqrt2\)。

	对每个 \(i\)，令 \(A_i=(ih,ih)\)、\(B_i=((i+1)h,ih)\)、\(C_i=((i+1)h,(i+1)h)\)。按 \((A_i,B_i,C_i)\) 取向的三角形为正向，且
	\[
	\partial[A_iB_iC_i]=[A_iB_i]+[B_iC_i]-[A_iC_i].
	\]
	将这些三角形的积分流相加得到 \(S_m\)。各三角形内部不交，故
	\[
	\partial S_m=P_m-T,\qquad
	\mathbf M(S_m)=m\frac{h^2}{2}=\frac1{2m}.
	\]
	在平坦范数定义中取零余项与填充 \(S_m\)，即得所需上界。于是 \(P_m\) 平坦收敛、从而弱收敛于 \(T\)，质量却始终为二，质量范数误差也不趋零。

	对最后一问，取紧支光滑函数 \(\phi\)，在 \(\operatorname{spt}P\) 的一个邻域内等于 \(x+y\)。边界公式给出
	\[
	P(\mathrm d\phi)=\partial P(\phi)=2.
	\]
	每个网格边的单位切向量是 \(\pm e_1\) 或 \(\pm e_2\)，而
	\(\mathrm d\phi=\mathrm dx+\mathrm dy\) 在这些方向上的绝对值都等于一。因此逐边估计便有
	\(2=|P(\mathrm d\phi)|\le\mathbf M(P)\)。这个估计使用网格允许的切向方向；它不声称 \(\mathrm dx+\mathrm dy\) 在所有单位方向上的模都小于一。网格流的质量下界二严格大于极限对角线的质量 \(\sqrt2\)，所以变形定理中的固定维数倍数不能仅靠细化这一坐标网格改成趋于一的倍数。
\end{solution}

\begin{exercise}\label{g3:ex:real-current-isoperimetric-amplitude}
	设 \(1\le k<n\)。令 \(Q=[0,1]^{k+1}\times\{0\}^{n-k-1}\) 取标准取向，\(Z=\partial[Q]\)，并对 \(0<\lambda<1\) 令 \(T_\lambda=\lambda Z\)。证明 \(T_\lambda\) 是紧支实系数正规闭流，且每个有限质量填充 \(V\) 若满足 \(\partial V=T_\lambda\)，都有
	\[
	\mathbf M(V)\ge c_{n,k}\lambda.
	\]
	由此说明，不存在对所有实系数正规闭 \(k\) 流统一成立的常数 \(C\)，使每个这样的流都可填充为
	\(\mathbf M(V)\le C\mathbf M(T)^{(k+1)/k}\)。指出整数流的证明在哪一步排除了这个反例。
\end{exercise}
\begin{solution}
	立方体的多面体边界是有限质量整数流，且 \(\partial Z=0\)。因此实数倍 \(T_\lambda\) 正规、紧支、闭，并且
	\[
	\mathbf M(T_\lambda)=\lambda\mathbf M(Z)=2(k+1)\lambda.
	\]
	取紧支光滑 \(k\) 形式 \(\alpha\)，在 \(Q\) 的邻域内等于
	\(x_1\,\mathrm dx_2\wedge\cdots\wedge\mathrm dx_{k+1}\)。Stokes 公式给出
	\[
	Z(\alpha)=[Q](\mathrm d\alpha)=1.
	\]
	记 \(L=\sup_x\|\mathrm d\alpha_x\|_*<\infty\)。对任何所述填充 \(V\)，由边界定义和质量对偶界，
	\[
	\lambda=T_\lambda(\alpha)=V(\mathrm d\alpha)
	\le L\mathbf M(V).
	\]
	所以可取 \(c_{n,k}=1/L>0\)，且不需要预先知道 \(V\) 的支集。

	若存在题设统一上界，则同一个填充必须满足
	\[
	c_{n,k}\lambda
	\le C\bigl(2(k+1)\lambda\bigr)^{(k+1)/k}.
	\]
	除以 \(\lambda\)，再令 \(\lambda\downarrow0\)，右侧趋零，矛盾。实系数流可以通过缩小重数任意减小质量，几何尺度却完全不变。

	整数流等周填充证明使用了：固定 \(h\) 网格上，非零整数多面体 \(k\) 流的质量至少为 \(h^k\)。实系数多面体流的系数可任意接近零，这一离散下界不成立，所以不能从质量上界迫使网格流消失。
\end{solution}

\begin{exercise}\label{g3:ex:current-nontangent-weak-limit}
	在 \(\mathbb R^2\) 中，对 \(j\ge1\) 令 \(T_j\) 是 \(j\) 条水平线段
	\[
		[0,1/j]\times\{i/j\},\qquad i=1,\ldots,j,
	\]
	从左向右取向后对应整数流之和。计算 \(\mathbf M(T_j)\) 与 \(\mathbf M(\partial T_j)\)，求 \(T_j\) 的弱极限 \(T\)。证明 \(\partial T\) 不具有局部有限质量，并说明即使 \(\|T\|\) 集中于一条直线，\(T\) 也不是沿该直线的可整流流。
\end{exercise}
\begin{solution}
	不同线段互不相交，每条长 \(1/j\)，所以
	\[
		\mathbf M(T_j)=1,\qquad
		\partial T_j=\sum_{i=1}^j
		\bigl(\delta_{(1/j,i/j)}-\delta_{(0,i/j)}\bigr),
		\qquad \mathbf M(\partial T_j)=2j.
	\]
	对紧支光滑一形式 \(\omega=a(x,y)\,\mathrm dx+b(x,y)\,\mathrm dy\)，有
	\[
		T_j(\omega)=\sum_{i=1}^j\int_0^{1/j}a(x,i/j)\,\mathrm dx.
	\]
	将每项中的 \(a(x,i/j)\) 换成 \(a(0,i/j)\)，总误差不超过
	\(\|\partial_xa\|_\infty/(2j)\)。剩下的是 Riemann 和，故
	\[
		T(\omega)=\int_0^1a(0,y)\,\mathrm dy,\qquad
		\|T\|=\mathcal H^1\llcorner(\{0\}\times[0,1]).
	\]
	这个向量测度在竖直线段上处处指向水平方向。

	对测试函数 \(\varphi\)，
	\[
		\partial T(\varphi)=T(\mathrm d\varphi)
		=\int_0^1\partial_x\varphi(0,y)\,\mathrm dy.
	\]
	取非负 \(\eta\in C_c^\infty((0,1))\)，使 \(\int\eta=1\)，再取 \(\chi\in C_c^\infty(\mathbb R)\)，在零的邻域为一。令
	\(\varphi_m(x,y)=\chi(x)\eta(y)\sin(mx)\)。
	这些函数支于一个共同紧集，上确界范数一致有界，而
	\[
		\partial T(\varphi_m)=m.
	\]
	因此 \(\partial T\) 不能由局部有限测度表示，\(T\) 不是正规流。

	若 \(T\) 是可整流流，其取向应在 \(\|T\|\)-几乎每点平行于竖直切线；于是它作用在 \(a\,\mathrm dx\) 上应为零。实际取 \(a\) 在该线段附近等于一时，\(T(a\,\mathrm dx)=1\)，矛盾。这个例子同时显示：有界质量不能控制极限边界，而测度支集的可整流性也不能代替向量的切向性。
\end{solution}

\begin{exercise}\label{g3:ex:current-tightness-global-flat}
	设 \(0\le k\le n\)，\(T_j\in\mathcal I_k(\mathbb R^n)\)，并且
	\[
		\sup_j\bigl(\mathbf M(T_j)+\mathbf M(\partial T_j)\bigr)<\infty,
		\qquad
		\lim_{R\to\infty}\sup_j
		\|T_j\|(\mathbb R^n\setminus B_R)=0.
	\]
	证明存在子列及 \(T\in\mathcal I_k(\mathbb R^n)\)，使
	\(\mathcal F(T_j-T)\to0\)。证明不必额外假设边界质量在无穷远处一致紧。解释为什么正文中平移单位圆的例子不满足上述条件。
\end{exercise}
\begin{solution}
	局部整数流紧性给出子列 \(T_j\rightharpoonup T\)，其中 \(T\) 局部整数。对递增球使用质量下半连续性并令半径趋于无穷，得到
	\[
		\mathbf M(T)+\mathbf M(\partial T)
		\le\liminf_j
		\bigl(\mathbf M(T_j)+\mathbf M(\partial T_j)\bigr)<\infty.
	\]
	故 \(T\) 是全空间整数流。

	给定 \(\varepsilon>0\)，选 \(R\) 使所有 \(j\) 都有
	\(\|T_j\|(\mathbb R^n\setminus B_R)<\varepsilon\)。由于
	\(\mathbb R^n\setminus\overline{B_R}\) 是开集，局部开集下半连续性再穷尽到无界开集，得到
	\[
		\|T\|(\mathbb R^n\setminus\overline{B_R})\le\varepsilon.
	\]
	取 \(0\le\eta\le1\) 光滑紧支，在 \(\overline{B_R}\) 的邻域为一。正文的局部平坦结论说明
	\(\mathcal F(\eta T_j-\eta T)\to0\)。
	因为 \(1-\eta\) 在 \(\overline{B_R}\) 的邻域为零，
	\[
		\mathbf M((1-\eta)T_j)\le\varepsilon,\qquad
		\mathbf M((1-\eta)T)\le\varepsilon.
	\]
	用 \(\mathcal F(A)\le\mathbf M(A)\) 与三角不等式，得到
	\[
		\limsup_j\mathcal F(T_j-T)
		\le 2\varepsilon.
	\]
	令 \(\varepsilon\downarrow0\)，便得所需收敛。这里边界的总质量界用于正规流紧性，而截去远处部分的平坦误差直接由该部分的质量控制，所以不需要再对边界尾部加条件。

	平移单位圆 \(C_j\) 的全部质量最终位于任意固定球外，故对每个 \(R\) 都有
	\(\sup_j\|C_j\|(\mathbb R^2\setminus B_R)=2\pi\)。它违反的正是所给质量紧性条件。
\end{solution}
